% Full manuscript: reworking-v2, 27 October 2026.
% Incorporates the revised Section 2 and Appendix C; numbered statements unchanged.
% Source: preamble
\documentclass[11pt,letterpaper]{article}
\usepackage[margin=1.05in]{geometry}
\usepackage[T1]{fontenc}
\usepackage{lmodern,microtype}
\usepackage{amsmath,amssymb,amsthm,mathtools}
\usepackage[hidelinks,hypertexnames=false,linktoc=all]{hyperref}
\newcommand{\Ac}{\mathsf A}
\usepackage{fancyhdr}
\pagestyle{fancy}
\fancyhf{}
\fancyhead[L]{\small Gaussian approximation over rectangles}
\fancyhead[R]{\small\thepage}
\setlength{\headheight}{14pt}
\hypersetup{pdftitle={High-Dimensional Gaussian Approximations over Rectangles without Covariance Restrictions},pdfsubject={Bounded independent summands; rates optimal up to logarithmic factors},pdfauthor={Victor Chernozhukov, Denis Chetverikov, Kengo Kato, Yuta Koike}}
\setlength{\parindent}{1.25em}
\setlength{\parskip}{0pt}
\emergencystretch=2em
\allowdisplaybreaks[2]
\raggedbottom
\clubpenalty=10000
\widowpenalty=10000
\newtheorem{theorem}{Theorem}
\newtheorem{proposition}{Proposition}
\newtheorem{lemma}{Lemma}
\newtheorem{corollary}{Corollary}
\theoremstyle{remark}
\newtheorem{remark}{Remark}
\newcommand{\E}{\mathbb E}
\newcommand{\Pp}{\mathbb P}
\newcommand{\R}{\mathbb R}
\newcommand{\cR}{\mathcal R}
\newcommand{\cF}{\mathcal F}
\newcommand{\cP}{\mathcal P}
\newcommand{\cD}{\mathcal D}
\newcommand{\Cov}{\operatorname{Cov}}
\newcommand{\Var}{\operatorname{Var}}
\newcommand{\tr}{\operatorname{tr}}
\newcommand{\op}{\mathrm{op}}
\newcommand{\norm}[1]{\left\lVert#1\right\rVert}
\newcommand{\rh}{\rho_{\mathcal R}}
\newcommand{\ind}{\mathbf 1}
\newcommand{\dd}{\,d}
\newcommand{\Hh}{H_2}
\newcommand{\surf}{\mathcal J}

\newcommand{\logp}{\log_+}

% Source: front
\begin{document}
\begin{center}
{\Large High-Dimensional Gaussian Approximations\\[0.15em]
over Rectangles without Covariance Restrictions}\\[0.85em]
{\normalsize Victor Chernozhukov, Denis Chetverikov, Kengo Kato, Yuta
Koike\renewcommand{\thefootnote}{\fnsymbol{footnote}}\footnote{The paper
was developed through collaboration with the AI systems GPT-6 Astra,
GPT-6 Sol, and Claude Opus.}}\\[0.65em]
{\normalsize 27 October 2026}
\end{center}
\setcounter{footnote}{0}
\thispagestyle{plain}

\begin{abstract}
 This paper studies the Gaussian approximation of sums of independent bounded random
vectors over rectangles, with no assumption on the correlation matrix,
which may be singular. Let \(X_1,\ldots,X_n\) be independent centered
random vectors in \(\R^p\) whose coordinates \(X_{ij}\) are bounded in
absolute value by a constant~\(B\), and suppose that each coordinate of
\(W_n=n^{-1/2}(X_1+\cdots+X_n)\) has unit variance. Let \(Z_n\) be a
centered Gaussian vector with the covariance matrix of \(W_n\), and put
\(L=\log(ep)\) and \(\log_+x=\max\{0,\log x\}\). The largest difference
between \(\Pp(W_n\in A)\) and \(\Pp(Z_n\in A)\) over rectangles \(A\) is
at most
\[
 \frac{BL^{3/2}}{\sqrt n}\sqrt{1+\log_+\sqrt{n/L}},
\]
up to a numerical constant. If, for a fixed \(c_0>0\), no two
coordinates have correlation larger than \(1-c_0\) in absolute value, the
bound improves slightly. When \(p\) grows
polynomially in~\(n\), the two bounds are of order \(n^{-1/2}(\log n)^2\) and
\(n^{-1/2}(\log n)^{3/2}\), respectively. A further refinement holds if the
third moments cancel in aggregate, that is, if
\(\sum_i\E[X_{ij}X_{ik}X_{i\ell}]=0\) for all coordinates \(j,k,\ell\).
The factor \(BL^{3/2}/\sqrt n\) cannot be improved in general, so the
second rate is then optimal and the first is optimal up to the factor
\((\log n)^{1/2}\).

\end{abstract}

\clearpage
\tableofcontents
\clearpage

\section{\texorpdfstring{Introduction and main results}{Introduction and main results}}\label{sec:results}
Gaussian approximation over rectangles provides a common foundation for
approximating maxima and constructing simultaneous confidence regions, with
applications to a wide variety of statistical inference problems reviewed in~\cite{CCKK2023}.
In many of these applications the correlation matrix may be nearly or
exactly singular, resulting in arbitrarily small eigenvalues. This motivates the
Gaussian approximations of this paper, whose quality does not depend on the
eigenvalues.

Consider a sum of independent centered bounded random vectors in
\(\R^p\) with marginal unit variance and covariance \(\Sigma\), and let
\(Z\sim N(0,\Sigma)\) be the ``approximating'' normal vector that shares the
mean and covariance of the sum:
\begin{equation}
 \begin{gathered}
 W=(W_j)_{j=1}^p=\sum_{i=1}^n\xi_i,\qquad
 \E\xi_i=0,\qquad \norm{\xi_i}_\infty\le\delta\quad\text{a.s.},\\
 \Sigma=\Cov(W),\qquad\Sigma_{jj}=1\ (j\le p),
 \qquad Z\sim N(0,\Sigma).
 \end{gathered}
 \label{eq:setup}
\end{equation}
The summands need not have the same distribution, and coordinates within
a summand may be dependent. The unit marginal variances fix the scale
of each coordinate; they can be obtained by coordinatewise rescaling
whenever all marginal variances are positive; there is no
restriction on the rank of \(\Sigma\).

The parameter \(\delta\) bounds the contribution of any one summand
to any coordinate of \(W\), and such a uniform bound is called an
\emph{envelope}. To state the
bound compactly, put
\[
 L=\log(ep),\qquad
 D_L=1+\log_+\frac1{\delta\sqrt L},\qquad
 \log_+x=\max\{0,\log x\}.
\]
Here \(L\) encodes the dependence on dimension, while \(D_L\)
contributes an additional logarithmic factor. Constants denoted by
\(C\) are absolute unless a subscript records their dependence.

\indent
Let \(\mathcal R\) denote the class of axis-aligned hyper-rectangles in
\(\R^p\). We measure the approximation error by the Kolmogorov distance
between the distributions of the random vectors \(W\) and \(Z\):
\[
 \rh(W,Z)=\sup_{A\in\mathcal R}
       |\Pp(W\in A)-\Pp(Z\in A)|.
\]
The same notation is used for any two random vectors in \(\R^p\). For
example, taking \(A=(-\infty,t]^p\) shows that this distance controls the
error in approximating the distribution of \(\max_{j\le p}W_j\), the maximum
coordinate of \(W\).

The first bound of this paper depends only on the summand envelope and the
dimension, so it holds uniformly as covariance eigenvalues approach or
reach zero.

\begin{theorem}[Unrestricted covariance]\label{thm:unrestricted}
Under~\eqref{eq:setup},
\begin{equation}
 \rh(W,Z)\le C\min\{1,\delta L^{3/2}\sqrt{D_L}\}.
 \label{eq:unrestricted}
\end{equation}
\end{theorem}

Theorem~\ref{thm:unrestricted} allows any rank of \(\Sigma\), perfect
correlations, and singular principal submatrices, and for normalized
observations it gives the following dependence on the sample size.

\begin{corollary}[Normalized observations]\label{cor:observations}
Let \(X_1,\ldots,X_n\in\R^p\) be independent and centered. Suppose
\(\norm{X_i}_\infty\le B\) almost surely and
\(n^{-1}\sum_i\E X_{ij}^2=1\) for every \(j\). Then, with
\(W_n=n^{-1/2}\sum_iX_i\) and \(Z_n\sim N(0,\Cov(W_n))\),
\begin{equation}
 \rh(W_n,Z_n)\le C\min\!\left\{1,
 \frac{B L^{3/2}}{\sqrt n}
 \sqrt{1+\log_+\frac{\sqrt n}{B\sqrt L}}\right\}.
 \label{eq:observations-unrestricted}
\end{equation}
\end{corollary}

Corollary~\ref{cor:observations} follows from Theorem~\ref{thm:unrestricted}
by setting \(\xi_i=X_i/\sqrt n\), which has envelope \(\delta=B/\sqrt n\). For
fixed \(B\) and \(p\) growing polynomially in~\(n\), its bound has
order \(n^{-1/2}(\log n)^2\), so the power \(n^{-1/2}\) of the
sample size is reached up to logarithmic factors, however small the
eigenvalues of \(\Sigma\) are.

The factor \(BL^{3/2}/\sqrt n\) is sharp. For \(B\ge2\), \(p\to\infty\) and
\(BL^{3/2}/\sqrt n\to0\), Theorem~\ref{opt:skew-lower} gives i.i.d.\ observations
with independent coordinates whose error is at least of order
\(BL^{3/2}/\sqrt n\). The example is inspired by that of Chernozhukov,
Chetverikov, and Koike~\cite[Proposition~2.1]{CCK}, and it drops some of their
side conditions on \(B\) and \(p\). Thus, in this regime, the
bound~\eqref{eq:observations-unrestricted} is optimal up to the logarithmic
factor \(\sqrt{1+\log_+(\sqrt n/(B\sqrt L))}\); when \(p\) grows polynomially
in~\(n\) and \(B\) is fixed, this factor is of order \((\log n)^{1/2}\). Whether
this factor can be improved remains an open question.

For unequal marginal variances, let
\[
 \sigma_j^2=n^{-1}\sum_i\E X_{ij}^2,
 \qquad
 \underline\sigma=\min_j\sigma_j>0.
\]
We then apply Theorem~\ref{thm:unrestricted} to the summands
\(\xi_{ij}=X_{ij}/(\sigma_j\sqrt n)\). If
\(\max_i\norm{X_i}_\infty\le K_n\), these summands have
envelope \(\delta=K_n/(\underline\sigma\sqrt n)\).
Positive diagonal scaling maps rectangles to rectangles and so preserves
the distance \(\rh\), and the resulting bound also holds in the
original coordinates.

\subsection{Correlation gaps and third-moment cancellation}
We can obtain a better bound when no two distinct coordinates are nearly
equal or opposite. To measure this separation, define
\[
 c=1-\max_{j\ne k}|\Sigma_{jk}|,
 \qquad c=1\ \text{when }p=1,
\]
and put
\begin{equation}
 \kappa=\begin{cases}\min\{L,c^{-1}\},&c>0,\\L,&c=0,\end{cases}
 \qquad E_\kappa=1+\log_+\frac{\sqrt\kappa}{\delta L}.
 \label{eq:parameters}
\end{equation}
A positive gap excludes equal or opposite coordinates; it still permits
singular covariance matrices.

\begin{theorem}[Gap-sensitive bound]\label{thm:optimized}
Under~\eqref{eq:setup},
\begin{equation}
 \rh(W,Z)\le C\min\!\left\{1,
 \delta[L^{3/2}+L\sqrt{\kappa E_\kappa}]\right\}.
 \tag{R*}\label{eq:optimized}
\end{equation}
\end{theorem}

Since \(\kappa\le L\) and \(E_\kappa\le D_L\),
Theorem~\ref{thm:optimized} implies Theorem~\ref{thm:unrestricted};
the term \(\delta L^{3/2}\) in~\eqref{eq:optimized} comes from the
third moments of the summands,
and it disappears when the full third-moment tensor cancels.

\begin{theorem}[Cancellation of third moments]\label{thm:zero-skew}
Under~\eqref{eq:setup}, suppose
\begin{equation}
 \sum_i\E\xi_i^{\otimes3}=0,
 \quad\text{equivalently,}\quad
 \sum_i\E[\xi_{ij}\xi_{ik}\xi_{i\ell}]=0
 \quad(j,k,\ell\le p).
 \label{eq:zero-third}
\end{equation}
Then
\begin{equation}
 \rh(W,Z)\le C\min\{1,\delta L\sqrt{\kappa E_\kappa}\}.
 \tag{ZS}\label{eq:zero-skew}
\end{equation}
\end{theorem}

Condition~\eqref{eq:zero-third} holds when each summand is centrally
symmetric, that is, when \(\xi_i\) and \(-\xi_i\) have the same law. Because it
involves the mixed third moments \(\E[\xi_{ij}\xi_{ik}\xi_{i\ell}]\), vanishing
marginal third moments do not suffice for this cancellation.

For a fixed standardized envelope and \(c\ge c_0>0\), the sample-size
bounds become
\begin{align}
 \rh(W_n,Z_n)&\le C_{B,c_0}\min\!\left\{1,
 \frac{L^{3/2}}{\sqrt n}
 +\frac L{\sqrt n}\sqrt{1+\log_+\frac{\sqrt n}{L}}\right\},
 \label{eq:observations-fixed}\\
 \rh(W_n,Z_n)&\le C_{B,c_0}\min\!\left\{1,
 \frac L{\sqrt n}\sqrt{1+\log_+\frac{\sqrt n}{L}}\right\}
 \quad\text{under~\eqref{eq:zero-third}}.
 \label{eq:zero-fixed}
\end{align}
In polynomial dimension the bound~\eqref{eq:observations-fixed}
has order \(n^{-1/2}(\log n)^{3/2}\), and it is optimal there, since the
lower-bound examples of~\cite[Proposition~2.1]{CCK} and
Theorem~\ref{opt:skew-lower} have independent coordinates and hence
correlation gap one. These bounds give uniform convergence at a fixed
positive gap when \(\log p=o(n^{1/3})\) in general and \(\log p=o(\sqrt n)\)
under cancellation. Theorem~\ref{opt:consistency} shows that these
conditions are sharp over fixed-envelope classes, with envelope at least
two for the first assertion and at least one for the second. In
the second case the sharpness concerns only the dimension threshold, and
below the threshold the error rate in~\eqref{eq:zero-fixed} is not known
to be sharp.

\subsection{Sub-exponential coordinates}\label{sec:subexponential}
The bounded theorems extend to sub-exponential coordinates if each
coordinate is truncated to an interval \([-T,T]\), at a
level \(T\) that depends only on \(\log n\) and on the tail scale \(K\) defined below, and the discarded sum is
controlled separately. We can measure the size of a random variable \(V\) through
the Orlicz norm
\[
 \|V\|_{\psi_1}=\inf\{t>0:\E e^{|V|/t}\le2\}.
\]
If the norm is finite, \(V\) is said to be sub-exponential.

\begin{corollary}[Sub-exponential observations]\label{cor:subexponential}
Let \(X_1,\ldots,X_n\in\R^p\) be independent and centered, with
\(n^{-1}\sum_i\E X_{ij}^2=1\) and
\(\max_{i,j}\|X_{ij}\|_{\psi_1}\le K\), where \(K\ge1\).
For \(W_n=n^{-1/2}\sum_iX_i\) and a centered Gaussian \(Z_n\)
with the same covariance, put
\[
 \ell_n=\log(en),\qquad b_n=\frac{K\ell_n}{\sqrt n},\qquad
 D_n=1+\log_+\frac1{b_n\sqrt L}.
\]
Then
\begin{equation}
 \rh(W_n,Z_n)\le C\min\{1,b_nL^{3/2}\sqrt{D_n}\}.
 \label{eq:subexponential}
\end{equation}
If the correlation gap of \(W_n\) is at least a fixed \(c_0>0\), then
\begin{equation}
 \rh(W_n,Z_n)\le C_{c_0}\min\!\left\{1,
 b_n\left[L^{3/2}+L\sqrt{1+\log_+\frac1{b_nL}}\right]\right\}.
 \label{eq:subexponential-gap}
\end{equation}
\end{corollary}

The bound~\eqref{eq:subexponential} requires neither a lower bound on a
local determinant nor a lower bound on the positive eigenvalues of the
covariance matrix of \(W_n\). For unequal marginal variances, the tail
condition is imposed on \(X_{ij}/\sigma_j\), as in the standardization after
Corollary~\ref{cor:observations}.

For fixed \(K\) and \(p\asymp n^a\) with \(a>0\) fixed, the
bound~\eqref{eq:subexponential} is of order \(n^{-1/2}(\log n)^3\), and the
bound~\eqref{eq:subexponential-gap} under a fixed correlation gap is of order
\(n^{-1/2}(\log n)^{5/2}\), so the power \(n^{-1/2}\) of the sample size is
again optimal up to logarithmic factors. By~\eqref{eq:subexponential}, the error
tends to zero for fixed \(K\) when \(\log p=o(n^{1/3}/\log(en))\).
The proof is in Appendix~\ref{app:subexponential}. It truncates each coordinate
at a level of order \(K\log(en)\). The exponential moments bound the discarded
part of the sum directly, so the level depends on \(K\) and \(\log n\) alone. Truncation
changes the covariance, and a joint Gaussian coupling of the truncated and the
discarded parts gives \(Z_n\) exactly the covariance of \(W_n\).

\subsection{Relation to recent work}\label{sec:literature}
High-dimensional Gaussian approximation for maxima and rectangles was
developed by Chernozhukov, Chetverikov, and Kato~\cite{CCK2013,CCK2017}.
Without covariance restrictions beyond positive marginal variances,
Chernozhukov, Chetverikov, Kato, and Koike~\cite[Theorem~2.1]{CCKK2022}
obtained an \(n^{-1/4}\) rate, up to logarithms, under coordinatewise
sub-exponential tails and an averaged fourth-moment condition.
The cube-root rectangle bound reported in~\cite[Theorem~1 and Section~5]{GaoZhang} improves the sample-size power to \(n^{-1/3}\). Under the tail assumption of Corollary~\ref{cor:subexponential}, in the unit-variance notation of this paper, it is
\begin{equation*}
 C\min\left\{1,
 \left(\frac{K^2}{n}\right)^{1/3}\{\log(2np)\}^{7/3}
 +\frac K{\sqrt n}\{\log(2np)\}^{5/2}\right\}.
\end{equation*}
No additional fourth-moment condition is needed for this bound.
For fixed \(K\) and \(p\asymp n^a\), with fixed \(a>0\), it has
order \(n^{-1/3}(\log n)^{7/3}\). In this dimension regime,
Corollaries~\ref{cor:observations} and~\ref{cor:subexponential}, with
fixed \(B\) or \(K\), give the sharp sample-size power
\(n^{-1/2}\), with logarithmic factors \((\log n)^2\) for bounded
observations and \((\log n)^3\) for sub-exponential coordinates.
These comparisons concern fixed envelopes or tail scales; the full
bounds retain their different dependence on these parameters.

A second line of work obtains rates close to \(n^{-1/2}\) under
covariance nondegeneracy. For bounded sums with positive definite
\(\Sigma\), the nondegenerate \(n^{-1/2}\) bound reported in~\cite[Theorem~1.1]{AL} removes the
\(\log(en)\) factor from the bound of Chernozhukov, Chetverikov,
and Koike~\cite[Corollary~2.1]{CCK}, and reads
\begin{equation*}
 \rh(W_n,Z_n)\le C\min\left\{1,
 \frac{BL^{3/2}}{\lambda_{\min}(\Sigma)\sqrt n}\right\}.
\end{equation*}
A fixed positive minimum eigenvalue therefore gives the sharp
\(BL^{3/2}/\sqrt n\) scale in the regimes of matching lower
bounds~\cite[Proposition~2.1]{CCK}.
The unrestricted bound of this paper does not involve \(\lambda_{\min}(\Sigma)\),
and it exceeds this scale only by the factor \(\sqrt{D_L}\).

Other results replace full nondegeneracy by conditions on small
groups of coordinates. Fang, Koike, Liu, and Zhao~\cite[Theorem~1.1]{FKLZ} obtain nearly
\(n^{-1/2}\) rates under sub-exponential tails and quantitative
two- and three-coordinate covariance conditions, allowing the full
covariance matrix to be singular. The companion paper~\cite{Local}
assumes uniformly positive eigenvalues for principal submatrices of
order at most four. For bounded observations it gives
\(C_aBL^{3/2}/\sqrt n\), and \(C_aBL/\sqrt n\) under aggregate
third-moment cancellation. The present results also cover singular
three- and four-coordinate submatrices. The two papers therefore
complement each other, and without the local condition the bound loses
only the factor \(\sqrt{D_L}\).

The bounds of this paper reach the rate \(n^{-1/2}\), up to logarithmic factors, with
no condition on the covariance matrix, and they lose only the factor
\(\sqrt{D_L}\) against the sharp bounds under nondegeneracy. They arise
from an interpolation between the sum and its Gaussian counterpart along
which the derivative of a smoothed rectangle probability is bounded
without reference to the eigenvalues of the covariance matrix.

Section~\ref{sec:smoothing} outlines the proof, derives
the rates from its main estimates, and explains why they reach the
power \(n^{-1/2}\), comparing the argument with those of \cite{AL} and~\cite{GaoZhang}; it also introduces the
notation used in the appendices. Section~\ref{opt:section}
discusses sharpness, and Subsection~\ref{sec:lean-certification}
records the formal verification in Lean.
Appendices~\ref{sec:main-proof}--\ref{app:subexponential}
contain all proofs of the upper bounds. The Supplement~\cite{Supplement} collects the lower-bound calculations,
exact coordinate reductions, and further investigations.

\subsection{\texorpdfstring{Formal verification in Lean}{Formal verification in Lean}}\label{sec:lean-certification}
\leavevmode Theorems~\ref{thm:unrestricted}--\ref{opt:consistency} and the two bounds of
Corollary~\ref{cor:subexponential} were certified on 20 and 21~October 2026
in Lean~4.19.0, a proof assistant that checks formal proofs by computer,
with its library mathlib~v4.19.0, and the proofs use only Lean's standard
axioms. The Lean sources, the axiom audits, and the reviews that compare the
formal statements with the manuscript at that time are in the
\href{https://github.com/VC2015/HDCLTs/blob/bbf1ab99c733a398c181ca8c9f5a243a41f258c6/verification/2026-10-21-companion-results/README.md}{verification record}
of the repository, for source revision
\href{https://github.com/VC2015/HDCLTs/commit/6548201a5c5d9c132a53f1d504f7ef70346aa782}{\texttt{6548201a5c5d}}.
The certification refers to these Lean proofs, not to the written
derivations, and the results of the Supplement are not certified.
The revised proofs in \texttt{reworking-v2} have not been re-certified.
Updating the Lean development and checking its correspondence with this
version of the manuscript remain to be done.

% Replace the complete Section 2 with this file; use the existing preamble.
% Insert together with Rectangle_CLT_Appendix_C-reworking-v2.tex.
% Prepared against Git blob 2076dea4e857bbd02277b4a6f5f92135ace8a065.
\section{Outline of the proof}
\label{sec:smoothing}

We replace the indicator of a rectangle by a logistic approximation and
compare its expectations along a path that gradually replaces the bounded
summands with Gaussian noise while preserving the covariance. We regard
the interpolation parameter as time: the path starts at the bounded sum
at time zero and ends at the Gaussian vector at time one. We estimate
the derivatives using the sum of the summands' covariance matrices.
The resulting bound does not depend on the eigenvalues or rank of this
matrix and therefore applies even when it is singular.

Differentiating these expectations with respect to time produces
derivatives of the logistic approximation, which we estimate in two
ways. At early times, we bound these derivatives directly. At later
times, when the path contains enough Gaussian noise, we use Gaussian
integration by parts. This rewrites expectations involving third- and
fourth-order derivatives of the logistic approximation in terms of its
first- and second-order derivatives, multiplied by explicit functions
of the Gaussian noise. The lower derivatives grow more slowly as the
approximation becomes sharper, but we must also control the additional
Gaussian factors. In both estimates, we also use anti-concentration
inequalities, which control the probability of falling in a narrow
region near the boundary of a rectangle. These probability bounds
provide an additional small factor. The more delicate part of the
proof is to preserve this gain when applying integration by parts.

Subsection~\ref{sec:outline-smoothing} introduces the logistic
approximation and bounds its smoothing error.
Subsection~\ref{sec:outline-path} introduces the calibrated
interpolation path and states the derivative bounds.
Subsections~\ref{sec:outline-early} and~\ref{sec:outline-late} explain
how to bound the derivatives directly at early times and how to use
Gaussian integration by parts at later times. In
Subsection~\ref{sec:rate}, we integrate these bounds and choose the
smoothing scale to balance the smoothing and comparison errors.

The proofs are organized by these dependencies. Appendix~\ref{sec:main-proof}
derives the theorems from the smoothing and derivative bounds.
Appendix~\ref{appx-sec:smoothing} proves the properties of the logistic
approximation, and Appendix~\ref{appx-sec:support} proves the derivative
bounds. Appendix~\ref{appx-app:concentration} gives the anti-concentration
argument, and Appendix~\ref{app:subexponential} treats sub-exponential
coordinates.

\subsection{\texorpdfstring{The logistic approximation}{The logistic approximation}}
\label{sec:outline-smoothing}

To construct this approximation, write the rectangle as
\(A=\prod_{j=1}^p[\ell_j,u_j]\), where \(\ell_j\le u_j\),
allowing infinite endpoints.
Each finite endpoint imposes one coordinate bound. If \(e_j\) denotes
the \(j\)th standard basis vector in \(\R^p\), the upper and lower
bounds can both be written with the same inequality sign:
\[
 x_j\le u_j
 \quad\Longleftrightarrow\quad e_j\cdot x\le u_j,
 \qquad
 x_j\ge\ell_j
 \quad\Longleftrightarrow\quad (-e_j)\cdot x\le-\ell_j.
\]
Here \(\cdot\) denotes the usual inner product on \(\R^p\).
Omit bounds with infinite endpoints and enumerate the remaining
inequalities by \(\alpha=1,\ldots,m\), where \(m\le2p\). We can
then write
\[
 A=\{x\in\R^p:\nu_\alpha\cdot x\le a_\alpha
                 \text{ for }\alpha=1,\ldots,m\}.
\]
We call \(\nu_\alpha\) the \emph{face normal} and \(a_\alpha\) the
\emph{threshold}. The face normal points outward: an upper bound uses
\((\nu_\alpha,a_\alpha)=(e_j,u_j)\), and a lower bound uses
\((\nu_\alpha,a_\alpha)=(-e_j,-\ell_j)\).
Thus the face normals are distinct members of
\(\{\pm e_1,\ldots,\pm e_p\}\).

For a smoothing scale \(h>0\), define the logistic approximation to
the indicator function of \(A\):
\begin{equation}
 F_{A,h}(x)=\left(1+\sum_{\alpha=1}^m
          e^{(\nu_\alpha\cdot x-a_\alpha)/h}\right)^{-1}.
 \label{eq:selector}
\end{equation}

For each fixed \(x\), this formula defines a probability distribution
over the \(m\) faces and an additional baseline option. Face
\(\alpha\) receives weight \(e^{(\nu_\alpha\cdot x-a_\alpha)/h}\),
while the baseline option receives weight one. After normalization,
the baseline probability is \(F_{A,h}(x)\), and the face probabilities
sum to \(1-F_{A,h}(x)\). We use these probabilities below to express
and bound the derivatives of the logistic approximation.

We compare the bounded sum \(W\) with the Gaussian vector \(Z\)
through the expectations of these smooth functions. Their largest
discrepancy is
\begin{equation}
 M_h=\sup_A\left|\E F_{A,h}(W)-\E F_{A,h}(Z)\right|,
 \label{eq:smoothed-discrepancy}
\end{equation}
where the supremum ranges over the rectangles described above.
The interpolation argument will provide a bound for \(M_h\).
Combining this bound with the error from replacing rectangle indicators
by their logistic approximations will then give the desired Gaussian
approximation bound.

Rectangle membership requires the largest coordinate-bound violation
to be nonpositive. The logistic approximation smooths both the maximum and the
indicator that tests this condition. For \(m\ge1\), we can write
\[
 F_{A,h}(x)=g_h\!\left(h\log\sum_{\alpha=1}^m
 e^{(\nu_\alpha\cdot x-a_\alpha)/h}\right),
 \qquad g_h(t)=\frac{1}{1+e^{t/h}}.
\]
The inner function smoothly approximates the maximum violation, while
\(g_h\) approximates the indicator of \(({-\infty},0]\). We use the
same scale \(h\) for both steps; the construction of Chernozhukov,
Chetverikov, and Kato~\cite{CCK2013} allows the cutoff function and
its scale to be chosen separately from the smooth maximum. The
resulting multinomial logistic form, with a baseline score of zero,
was used by McFadden~\cite[p.~1001]{McFadden1989} to smooth indicators
of simultaneous inequalities and also appears in~\cite[eq.~(5)]{GaoZhang}. For every \(h>0\), the approximation
is smooth, allowing us to differentiate along the interpolation path.
As \(h\) decreases to zero, it approaches the rectangle indicator
away from the boundary, but the uniform bounds on its derivatives
increase.

An anti-concentration inequality bounds the probability that a random
vector falls in a boundary layer, that is, within a given distance of
the boundary of a set. The following one controls both the error of
the logistic approximation and its derivatives near the boundary.
For \(A=\prod_j[\ell_j,u_j]\), write
\(A^t=\prod_j[\ell_j-t,u_j+t]\), with reversed intervals empty.
\begin{lemma}[Rectangle anti-concentration]\label{lem:concentration}
Let \(\zeta_1,\ldots,\zeta_n\) be independent centered random vectors
in \(\R^p\), and set \(S=\sum_{i=1}^n\zeta_i\). Suppose that
\[
 \norm{\zeta_i}_\infty\le\eta
 \quad\text{almost surely for every }i,
 \qquad
 \Var(S_j)=v>0\quad\text{for every }j.
\]
Let \(R\) be a random vector in \(\R^p\), independent of
\(\zeta_1,\ldots,\zeta_n\). Then, for every rectangle \(A\) and
\(t>0\),
\[
 \Pp(S+R\in A^t\setminus A^{-t})
 \le C\min\{1,(t+\eta)\sqrt{\log(ep)/v}\},
\]
where \(C\) is an absolute constant. If instead \(S\) is a centered
Gaussian vector with coordinate variances \(v>0\), the same bound
holds with \(\eta=0\) for every \(R\) independent of \(S\), including
when the covariance matrix of \(S\) is singular.
\end{lemma}
Appendix~\ref{appx-app:concentration} proves the bounded-sum
assertion by conditioning on pairs of independent copies and applying
the Rademacher anti-concentration inequality of O'Donnell, Servedio,
and Tan~\cite{OST}, in the form recorded
in~\cite[Lemma~J.4]{CCKK2022}. This recovers the bounded-sum estimate
in~\cite[Lemma~2.1]{AL}, obtained using
a bootstrap argument and Kane's~\cite{Kane} noise-sensitivity inequality.
The Gaussian assertion is Nazarov's inequality;
see~\cite[Lemma~A.1]{CCK2017}.

Let \(\ind_A\) denote the indicator function of \(A\): it equals one
on \(A\) and zero outside \(A\). Recall that
\(\kappa=\min\{L,c^{-1}\}\), interpreted as \(\kappa=L\) when
\(c=0\), where \(L=\log(ep)\) and \(c\) is the correlation gap.
The following lemma bounds the error from replacing the indicator by
its logistic approximation and relates \(M_h\) to the Kolmogorov
distance over rectangles.

\begin{lemma}[Smoothing error]\label{lem:endpoint}
Under~\eqref{eq:setup}, for \(h\ge\delta\),
\begin{equation}
 \E|F_{A,h}(W)-\ind_A(W)|+
 \E|F_{A,h}(Z)-\ind_A(Z)|\le Ch\sqrt L\,\kappa.
 \label{eq:endpoint-error}
\end{equation}
Consequently,
\begin{equation}
 \rh(W,Z)\le M_h+Ch\sqrt L\,\kappa.
 \label{eq:endpoint-min}
\end{equation}
\end{lemma}

To bound the smoothing error without assuming a positive correlation
gap, let \(\omega_0,\omega_1,\ldots,\omega_m\) be independent
exponential random variables with mean one. The logistic approximation
has the exact representation
\begin{equation}
 \begin{split}
 F_{A,h}(x)&=\Pp_\omega(x\in A_\omega),\\
 A_\omega&=\{x:\nu_\alpha\cdot x\le
 a_\alpha+h\log(\omega_\alpha/\omega_0)
 \text{ for every }\alpha\}.
 \end{split}
 \label{eq:mixture}
\end{equation}
Thus \(A_\omega\) is a rectangle with randomly shifted thresholds.
Let
\[
 R_h=\max_{\alpha\le m}
       |h\log(\omega_\alpha/\omega_0)|
\]
be the largest absolute change in a threshold, with \(R_h=0\) when
\(m=0\). Each normalized change has tails bounded by \(2e^{-t}\),
so a union bound over the \(m\le2p\) faces and integration give
\(\E R_h\le ChL\). For fixed exponentials, the indicators of \(A\)
and \(A_\omega\) can differ only in the boundary layer
\(A^{R_h}\setminus A^{-R_h}\). Applying
Lemma~\ref{lem:concentration} and averaging over the exponentials
therefore bounds the smoothing error by
\(C(\E R_h+\delta)\sqrt L\le ChL^{3/2}\), since \(h\ge\delta\).
When \(c>0\), we obtain another bound using the second derivatives of
\[
 \Psi(x)=h\log\!\left(1+\sum_\alpha
              e^{(\nu_\alpha\cdot x-a_\alpha)/h}\right).
\]
The pointwise error \(|F_{A,h}(x)-\ind_A(x)|\) is at most
\(2h\,\tr D^2\Psi(x)\), where the trace of the Hessian is the sum
of the coordinate second derivatives.
Lemma~\ref{appx-lem:endpoint-gap} bounds the expected trace at \(W\)
and \(Z\) by \(C\sqrt L/c\), giving the alternative smoothing
error \(Ch\sqrt L/c\). Taking the smaller of the two bounds gives
the factor \(\kappa\) in Lemma~\ref{lem:endpoint}.
Lemmas~\ref{appx-lem:endpoint-free} and~\ref{appx-lem:endpoint-gap}
provide the details of these two arguments.

The comparison along the interpolation path defined below in
Subsection~\ref{sec:outline-path} uses third- and fourth-order
derivatives of \(F_{A,h}\). Their size is controlled by the logistic
transition weight \(F(1-F)\), whose average we bound below. We
abbreviate \(F_{A,h}\) to \(F\) when \(A\) and \(h\) are fixed,
and write
\[
 w_\alpha(x)=e^{(\nu_\alpha\cdot x-a_\alpha)/h},\qquad
 \pi_\alpha=\frac{w_\alpha}{1+\sum_\beta w_\beta},\qquad
 \theta=1-F=\sum_\alpha\pi_\alpha,
\]
so that \(\pi_\alpha\) is the multinomial logistic probability of
face~\(\alpha\) and \(\theta\) is the total probability of the faces.

For the first derivative, the chain rule gives
\[
 \partial_jF(x)=-\frac{F(x)}h
       \sum_\alpha\pi_\alpha(x)\nu_{\alpha,j},
\]
where \(\nu_{\alpha,j}\) is the \(j\)th coordinate of the face
normal \(\nu_\alpha\). Each face normal is a signed coordinate
vector, so \(\sum_j|\nu_{\alpha,j}|=1\). Summing the absolute
values therefore gives
\[
 \sum_{j=1}^p|\partial_jF(x)|
 \le\frac{F(x)}h\sum_\alpha\pi_\alpha(x)
 =\frac{F(x)\theta(x)}h.
\]

Repeated differentiation introduces one factor \(h^{-1}\) at each
step. The remaining terms are products of \(F\), the face
probabilities, and coordinate components of the face normals. The
same reasoning controls their total absolute value: summing over
coordinates uses \(\sum_j|\nu_{\alpha,j}|=1\), while summing
products of \(k\) face probabilities gives \(\theta^k\le\theta\).
The number and coefficients of these terms depend only on the
derivative order.

We also need a bound that holds uniformly at nearby points. If
\(\norm z_\infty\le h\), each exponential weight
\(w_\alpha(x+z)\) differs from \(w_\alpha(x)\) by a factor
between \(e^{-1}\) and \(e\). Comparing each product at \(x+z\)
with its value at \(x\) therefore controls the derivatives
throughout this neighborhood. For an ordered tuple
\(I=(j_1,\ldots,j_q)\), write
\(\partial_I=\partial_{j_1}\cdots\partial_{j_q}\). We obtain the
\emph{derivative envelope}
\[
 \sum_{I\in\{1,\ldots,p\}^q}
 \sup_{\norm z_\infty\le h}|\partial_IF(x+z)|
 \le C_qh^{-q}F(x)\theta(x),\qquad q\ge1,
\]
where \(C_q\) depends only on \(q\).
Appendix~\ref{appx-sec:smoothing} proves this bound
as~\eqref{appx-eq:local-envelope}, using the derivative formula
in~\cite[Lemma~13]{GaoZhang}.

The product \(F\theta\) is at most \(1/4\), but anti-concentration
gives a smaller bound for its expectation. To see how, fix \(x\)
and write \(T=\sum_\alpha w_\alpha(x)\). Shifting every threshold
outward by \(th\) replaces \(T\) by \(e^{-t}T\). Differentiating
with respect to \(t\) and integrating therefore gives
\[
 F_{A^h,h}(x)-F_{A^{-h},h}(x)
 =\int_{-1}^{1}\frac{T}{(e^{t/2}+Te^{-t/2})^2}\,dt
 \ge\frac2e F(x)\theta(x),
\]
because both exponential factors in the denominator are at most
\(e^{1/2}\) throughout the integration interval.

By the random-threshold representation, the difference on the left
is an average of indicators of boundary layers of width \(2h\)
around randomly shifted rectangles. Applying
Lemma~\ref{lem:concentration} to these boundary layers gives
\[
 \E[F(U)\theta(U)]\le Ch\sqrt L,
\]
whenever \(U\) has an independent component satisfying the
assumptions of Lemma~\ref{lem:concentration} with coordinate variance
\(v\ge1/2\) and summand bound \(\eta\le h/2\).
Lemma~\ref{appx-lem:envelope} states the precise conditions.
Combining this estimate with the derivative envelope yields
\[
 \E\!\left[\sum_{I\in\{1,\ldots,p\}^q}
 \sup_{\norm z_\infty\le h}|\partial_IF(U+z)|\right]
 \le C_qh^{1-q}\sqrt L.
\]
Thus averaging supplies a factor \(h\sqrt L\) beyond the uniform
derivative bound \(C_qh^{-q}\). This gain is useful when
\(h\sqrt L\) is small.

\subsection{The Bernoulli--Gauss interpolation}
\label{sec:outline-path}

The Bernoulli--Gauss interpolation of~\cite[Section~4]{AL} removes
summands at random while increasing the Gaussian variance. We use
the same thinning probability and amplitude, with logistic smoothing
at the bounded endpoint. Let \(\lambda\in[0,1]\) denote the
interpolation time, let \(Z\) be independent of the summands, and set
\begin{equation}
 \gamma=1+\lambda,\qquad
 \varpi=\frac{1-\lambda}{(1+\lambda)^2},\qquad
 U_\lambda=\gamma\sum_i\mathbb B_{i,\lambda}\xi_i+\sqrt\lambda Z,
 \quad\mathbb B_{i,\lambda}\sim\operatorname{Bernoulli}(\varpi),
 \label{eq:path}
\end{equation}
with independent masks, summands, and Gaussian vector. Each mask
either retains or deletes its summand. The amplitude \(\gamma\)
ensures
\[
 \varpi\gamma^2=1-\lambda,\qquad
 \varpi\gamma^3=1-\lambda^2.
\]
The first identity keeps the total covariance equal to \(\Sigma\),
with \(U_0\stackrel d=W\) and \(U_1\stackrel d=Z\).
The second makes the derivative of the cubic moment coefficient
equal to \(-2\lambda\), rather than a nonzero constant at zero.
This calibration is the reason for the improved behavior near the
bounded endpoint~\cite[eq.~(17)]{AL}.

Put \(P_A(\lambda)=\E F_{A,h}(U_\lambda)\). The smoothing lemma
reduces the theorem to the integrated derivative:
\begin{equation}
 \rh(W,Z)\le M_h+Ch\sqrt L\,\kappa,
 \qquad M_h\le\sup_A\int_0^1|P_A'(\lambda)|\dd\lambda.
 \label{eq:roadmap}
\end{equation}
Differentiating one summand law at a time creates a comparison between
the full path and a path with that summand deleted. The Gaussian heat
derivative cancels the quadratic contribution because their total
covariances agree. The signed Hessian identity
\eqref{appx-an:signed-hessian} makes this cancellation explicit; its
first two moments then give
\begin{equation}
 P_A'(\lambda)=-\frac\lambda3
       \E D^3F_{A,h}(U_\lambda):\sum_i\E\xi_i^{\otimes3}
       +R_{4,A,h}(\lambda).
 \label{eq:path-summary}
\end{equation}
The colon contracts all three indices. Lemma~\ref{appx-an:lem:identity}
gives the exact fourth-order remainder, including the cost of moving
the cubic term from the deleted-summand bases to the full path.
Every cubic tensor is therefore evaluated against the same derivative
tensor, and under~\eqref{eq:zero-third} the cancellation occurs before
any absolute value is taken:
\begin{equation}
 P_A'(\lambda)=R_{4,A,h}(\lambda).
 \label{zs:identity}
\end{equation}
Each remainder term is a fourth derivative at a translate
\(U_\lambda+v_i\), with \(\norm{v_i}_\infty\le4\delta\).
Its directions are fresh copies of summand \(i\), independent of the
full path, and two of the directions are the same copy. That repeated
pair is what lets the covariance identity control the sum over \(i\).

\begin{proposition}[Two time ranges]\label{an:prop:derivative}
There are sufficiently large absolute constants \(\Ac,d_0\) such that,
for \(h\ge\Ac\delta\), uniformly over the thresholds,
\begin{align}
 |P_A'(\lambda)|&\le C\sqrt L\left(
 \frac{\lambda\delta}{h^2}+\frac{\delta^2}{h^3}\right),
 &&0<\lambda<1,\label{an:early}\\
 |P_A'(\lambda)|&\le CL^{3/2}\left(
 \delta+\frac{\delta^2}{\lambda h}\right),
 &&d_0\delta^2L\le\lambda<1.\label{an:late}
\end{align}
The fourth-order remainder separately satisfies
\begin{align}
 |R_{4,A,h}(\lambda)|&\le C\delta^2\sqrt L/h^3,
 &&0<\lambda<1,\label{eq:remainder-early}\\
 |R_{4,A,h}(\lambda)|&\le C\delta^2L^{3/2}/(\lambda h),
 &&d_0\delta^2L\le\lambda<1.\label{an:remainder}
\end{align}
\end{proposition}

Boundedness and continuity of the logistic approximation, together
with the integrable early bound, justify integration to both endpoints.

\subsection{The early bounds and the covariance identity}
\label{sec:outline-early}

We first isolate the covariance property used in both time ranges.
Factor \(\Sigma=\Gamma\Gamma^\top\), with \(\Gamma\) of full column
rank \(r\). The summands lie in its range, so
\(\xi_i=\Gamma\zeta_i^0\), and the face normals become
unit vectors \(u_\alpha=\Gamma^\top\nu_\alpha\). In these coordinates,
\[
 \sum_i\E\zeta_i^0(\zeta_i^0)^\top=I_r,
 \qquad \max_\alpha|u_\alpha\cdot\zeta_i^0|\le\delta.
\]
The face projections remain bounded even when the positive eigenvalues
of \(\Sigma\) are small. For fresh copies \(\zeta_i\), the covariance
identity and Cauchy--Schwarz imply, for any unit vectors \(u,v\),
\[
 \sum_i\E|(u\cdot\zeta_i)(v\cdot\zeta_i)|\le1.
\]
We keep the notation \(F(x)=F_{A,h}(\Gamma x)\) for the restricted
logistic function. Appendix~\ref{appx-sec:support} gives the coordinate
justification and the conditional version of this inequality.

In each third- or fourth-order derivative, use this inequality for
the repeated pair and the bound \(\delta\) for the other directions.
The resulting moment factors are \(\delta\) and \(\delta^2\).
The local derivative envelope absorbs shifts of size \(4\delta\le h\),
leaving the derivatives at the common path; their positive weights
are then independent of the fresh copies. Concentration supplies
\(\E F(U_\lambda)\theta(U_\lambda)\le Ch\sqrt L\), in support
coordinates, uniformly over \(\lambda\). Derivatives of order three
and four consequently contribute
\(C\delta h^{-2}\sqrt L\) and \(C\delta^2h^{-3}\sqrt L\), respectively.
The additional coefficient \(\lambda\) of the cubic term gives
\eqref{an:early} and~\eqref{eq:remainder-early}.

\subsection{The late bounds}
\label{sec:outline-late}

Once \(\lambda\ge d_0\delta^2L\), we can estimate the same derivatives
more effectively by moving the repeated pair onto the Gaussian density.
We use a single decomposition for every \(\lambda\). In support
coordinates, let \(V_\lambda\) be the bounded part of the path and write
\[
 U_\lambda=V_{0,\lambda}+\sigma_\lambda G_1,\qquad
 V_{0,\lambda}=V_\lambda+\sqrt{\lambda/2}\,G_0,\qquad
 \sigma_\lambda^2=\lambda/2,
\]
where \(G_0,G_1\) are independent standard Gaussian vectors.
The two independent components of \(V_{0,\lambda}\) have equal face
variances \(1-\lambda\) and \(\lambda/2\). At least one is at least
\(1/3\), so Lemma~\ref{lem:concentration} supplies concentration
uniformly along the path. We call such a translation a concentrated
base; its precise definition is~\eqref{appx-an:concentration}.

Write \(U=V_0+\sigma G\) for this decomposition. Conditional on the
base and the fresh directions, two integrations by parts give
\[
 \E_G D^{s+2}F(U)[\zeta_i,\zeta_i,z_1,\ldots,z_s]
 =\sigma^{-2}\E_G\left[
 \tr\{(GG^\top-I_r)\zeta_i\zeta_i^\top\}
 D^sF(U)[z_1,\ldots,z_s]\right],\qquad s=1,2.
\]
The remaining derivatives contain products of logistic weights and
\(s\) bounded face projections. At translated bases the weights
still depend on \(i\). The exact translation formula
\eqref{eq:weight-translation} expresses them as an absolutely
convergent series of weights at \(U\), with all dependence on the
summands in their coefficients. Only first and second derivatives
are needed after the integrations by parts; their explicit formulas
appear in Lemma~\ref{appx-an:lem:restoration}.

We then sum the coefficient matrices using the conditional covariance
bound. It controls their operator norms, and the inequality
\(|\tr(\Xi M)|\le\norm\Xi_\op\norm M_*\) reduces the problem to
Gaussian matrix averages. Here \(\norm M_*\) is the nuclear norm,
the sum of the absolute eigenvalues for a symmetric matrix. Summing
before taking this norm avoids a factor equal to the rank.
Lemma~\ref{appx-an:lem:matrix} proves
\[
 \E_{V_0}\sum_{\alpha_1,\ldots,\alpha_k}
 \left\|\E_G\left[(GG^\top-I_r)F(U)
                 \prod_{a=1}^k\pi_{\alpha_a}(U)\right]\right\|_*
 \le C^{k+1}hL^{3/2},\qquad \delta\sqrt L\le\sigma\le1.
\]
Consequently, integration by parts, the remaining directions, and
the matrix estimate contribute, respectively,
\[
 \sigma^{-2}\cdot\delta^sh^{-s}\cdot hL^{3/2}
 =\sigma^{-2}\delta^sh^{1-s}L^{3/2}.
\]
For \(s=1\), the coefficient \(\lambda\) of the cubic term cancels
\(\sigma^{-2}=2/\lambda\), giving \(C\delta L^{3/2}\).
For \(s=2\), the fourth-order terms give
\(C\delta^2L^{3/2}/(\lambda h)\). These are the late bounds.

The small factor in the matrix estimate deserves explanation.
A Gaussian law tilted by a log-concave logistic weight has covariance
at most \(I_r\). Its matrix contribution is therefore bounded by
a tilted-mean term and a signed trace term. The mean term inherits
the small total weight \(\E F\theta^k(U)\le Ch\sqrt L\).
The trace terms must first be summed; their sum becomes \(h/\sigma\)
times an averaged Gaussian surface integral. Lemma~\ref{appx-an:lem:surface}
provides the compensating factor \(\sigma\):
\[
 \E\left|\surf((Q-V_0)/\sigma)\right|
 \le C(\sigma L^{3/2}+\delta L^2)
 \le C\sigma L^{3/2},\qquad \sigma\ge\delta\sqrt L,
\]
where \(\surf(P)=\int_{\partial P}(\norm g^2-r)\varphi_r(g)
\dd\mathcal H^{r-1}(g)\). Its proof pairs the contributions from
adjacent facets, then uses outward tubes with widths chosen from
one- and two-dimensional normal projections. Let \(R(G)\) be one
plus the largest norm of these projections. After conditioning on
\(G\), concentration controls a boundary layer of width
\(\sigma/R(G)\); the second and third moments of \(R(G)\), of
orders \(L\) and \(L^{3/2}\), complete the bound.

% Source: rate_explanation
\subsection{\texorpdfstring{Integrating the comparison and choosing the scale}{Integrating the comparison and choosing the scale}}\label{sec:rate}
The early and late fourth-order bounds,
\(C\delta^2\sqrt L/h^3\) and \(C\delta^2L^{3/2}/(\lambda h)\), of
Proposition~\ref{an:prop:derivative} agree at the time
\(\tau_h=Lh^2\). Integrating the early bounds before this time and the
late bounds after it gives the following bound for the largest
difference \(M_h=\sup_A|\E F_{A,h}(W)-\E F_{A,h}(Z)|\) of smoothed
expectations, defined in~\eqref{eq:smoothed-discrepancy}, since
\(M_h\le\sup_A\int_0^1|P_A'(\lambda)|\dd\lambda\)
by~\eqref{eq:roadmap}. Its proof is in Appendix~\ref{sec:main-proof}.
We call a finite system of coordinate inequalities
\(\nu_\alpha\cdot x\le a_\alpha\) a threshold system, and call it
inconsistent if its inequalities have no common solution; such systems
arise in the proofs when thresholds are shifted.

\begin{proposition}[Smoothed comparison]\label{prop:moving-cutoff}
\label{zs:smooth-proposition}
There are absolute constants \(\mathsf A,C\) such that, under
\eqref{eq:setup}, for every \(h\ge\mathsf A\delta\),
\begin{equation}
 M_h\le CL^{3/2}\left[
 \delta+\frac{\delta^2}{h}\log_+\frac1{Lh^2}\right].
 \tag{S*}\label{eq:smooth-moving}
\end{equation}
If, in addition, the aggregate third-moment tensor vanishes,
as in~\eqref{eq:zero-third}, then
\begin{equation}
 M_h\le\frac{C\delta^2L^{3/2}}h
 \left[1+\log_+\frac1{Lh^2}\right].
 \label{zs:smooth}
\end{equation}
Both estimates are uniform over all finite threshold systems formed from
coordinate inequalities, including inconsistent systems.
\end{proposition}

In~\eqref{eq:smooth-moving}, the cubic term contributes the constant
\(C\delta L^{3/2}\) at every time, while the factor \(1/\lambda\) in the
late fourth-order bound produces the logarithm. Under cancellation the
cubic term is absent, so the early fourth-order contribution, which
\eqref{eq:smooth-moving} absorbs into \(C\delta L^{3/2}\), must be kept
explicitly; it accounts for the additive one in~\eqref{zs:smooth}.

Lemma~\ref{lem:endpoint} bounds the smoothing error by
\(Ch\sqrt L\,\kappa\), and balancing this against the comparison
remainder leads to the scale
\begin{equation}
 h=\Ac\delta\sqrt{LE_\kappa/\kappa}.
 \label{eq:optimal-scale}
\end{equation}
Here \(\kappa=\min\{L,c^{-1}\}\), or \(\kappa=L\) when \(c=0\), and
\(E_\kappa=1+\log_+\{\sqrt\kappa/(\delta L)\}\) are the parameters
of~\eqref{eq:parameters}, where \(c\) is the correlation gap.
Since \(\kappa\le L\) and \(E_\kappa\ge1\), this scale is at
least \(\Ac\delta\). Also
\(1+\log_+(1/(Lh^2))\le CE_\kappa\), so both the smoothing
error and the fourth-order comparison error are at most
\(C\delta L\sqrt{\kappa E_\kappa}\).
This yields Theorems~\ref{thm:optimized} and~\ref{thm:zero-skew};
Theorem~\ref{thm:unrestricted} follows from \(\kappa\le L\) and
\(E_\kappa\le D_L\). Appendix~\ref{sec:main-proof} gives the complete proof, including
the case \(Lh^2\ge1\), in which the late range is empty, and the
admissibility of the smoothing scale.

\subsection{\texorpdfstring{Sources of the argument and the power \(n^{-1/2}\)}{Sources of the argument and the power n to the -1/2}}
\label{sec:why-rate}\label{sec:sources}

The Bernoulli--Gauss path and its moment calibration are taken
from~\cite[Section~4]{AL}, as is the bounded-sum concentration
inequality~\cite[Lemma~2.1]{AL}. The fourth-order kernel in
Lemma~\ref{appx-an:lem:identity} is the kernel
of~\cite[Lemma~4.1]{AL}. The signed Hessian representation makes
its two cancellations explicit, while moving the cubic term to the
full path gives the same structure as the two-stage derivative
identity in~\cite[Lemma~22]{GaoZhang}. Keeping that common base
allows aggregate third moments to cancel exactly.

The support-coordinate reduction follows~\cite[Section~2.3]{GaoZhang}.
The late argument uses the covariance summation and two Gaussian
integrations by parts of~\cite[Lemma~23]{GaoZhang}; the exact
translation series restores a common base without a truncation
remainder. The matrix and surface estimates adapt the convex-cell
and polyhedral arguments of~\cite[Lemma~12 and Proposition~4]{GaoZhang},
including Gaussian covariance contraction and the pairing of the
boundary terms on each ridge. The multinomial logistic approximation
appears in~\cite[p.~1001]{McFadden1989} and~\cite[eq.~(5)]{GaoZhang};
its derivative and mixture formulas correspond to
\cite[Lemmas~13 and~14]{GaoZhang}.

The power of the sample size is decided by the smoothing error and
the integrated fourth-order remainder. In the notation used here,
the first-stage remainder estimate in~\cite[Proposition~2]{GaoZhang}
has order \(\delta^2L/(\lambda h^2)\). Its integral has a factor
\(h^{-2}\); balancing this against an error proportional to \(h\)
gives \(h\) of order \(\delta^{2/3}\), up to logarithmic factors,
and hence the power \(n^{-1/3}\) when \(\delta=B/\sqrt n\).
The late remainder here has order
\(\delta^2L^{3/2}/(\lambda h)\), smaller by \(h\sqrt L\).
Its integral has only \(h^{-1}\), so the balance uses \(h\) of
order \(\delta\), up to logarithmic factors, and gives \(n^{-1/2}\).
The bound in~\cite{AL} also reaches that power, but its Gaussian
derivative estimates depend on negative powers of
\(\lambda_{\min}(\Sigma)\).

The factor \(h\sqrt L\) comes from concentration applied to the
randomly shifted boundary layers in the logistic mixture. It enters
the early derivative bound directly. For the late bound, it must
survive the integrations by parts. Applying the matrix estimate to
the logistic weights themselves retains their small total mass.
The remaining trace term is multiplied by \(h/\sigma\), and the
averaged surface estimate supplies \(\sigma\), leaving \(h\).
A surface bound uniform over fixed translations would not supply
this factor. Averaging over the concentrated base is therefore
essential to the late bound and to the final rate.

% Source: sharpness_and_discussion
\section{Sharpness and remaining questions}\label{opt:section}
The dependence on sample size cannot improve beyond \(n^{-1/2}\)
under the present assumptions. A normalized Rademacher sum has an atom
of order \(n^{-1/2}\), while the Gaussian probability of an interval
shrinking to this atom tends to zero. This obstruction already occurs
in one dimension and survives zero skewness, that is, the
cancellation~\eqref{eq:zero-third} of third moments.

Dimension introduces another obstruction, already visible when the
coordinates are independent. Chernozhukov, Chetverikov, and
Koike~\cite[Proposition~2.1]{CCK} showed with i.i.d.\ two-point
coordinates that the factor \((\log p)^{3/2}\) in the rate cannot be
improved, under further growth conditions on \(B_n\) and \(p_n\) that the
following theorem does not need.
\begin{theorem}[Lower scale]\label{opt:skew-lower}
Suppose \(p_n\to\infty\), \(B_n\ge2\), and
\(r_n=B_n(\log p_n)^{3/2}/\sqrt n\to0\). There are i.i.d.
centered observations with independent coordinates, covariance
\(I_{p_n}\), and envelope \(B_n\), for which
\begin{equation}
 \liminf_n\frac{\rh(W_n,Z_n)}{r_n}\ge\frac{\sqrt2}{4e}.
 \label{opt:lower-constant}
\end{equation}
Here \(W_n=n^{-1/2}\sum_iX_{n,i}\),
\(Z_n\sim N(0,I_{p_n})\), and the envelope is attained.
\end{theorem}
The example uses the centered, unit-variance two-point law
\[
 \Pp(X=B_n)=\frac1{B_n^2+1},\qquad
 \Pp(X=-B_n^{-1})=\frac{B_n^2}{B_n^2+1}.
\]
At a normal maximum threshold of order \(\sqrt{\log p_n}\), its
skewness changes each relative tail by order \(r_n\). Taking the
product of the coordinate probabilities preserves this discrepancy.
The Supplement~\cite[Section~S.4.1, Theorem~S.3 and Lemma~S.2]{Supplement}
proves a moderate-tail estimate uniform in the envelope and includes the
lattice corrections. This extends the regime of the two-point lower
construction in~\cite[Proposition~2.1]{CCK}, with the law parameterized here
by its attained envelope.

These lower examples have correlation gap one, so they belong to
every class with a fixed positive gap. To express uniformity,
let \(M_n(c_0;B,p)\) be the supremum of rectangle error over the
independent centered observations with envelope at most \(B\), unit
average marginal variances, and correlation gap at least \(c_0\).
Write \(M_n^0(c_0;B,p)\) for the same supremum restricted to
observations whose aggregate third-moment tensor vanishes.
For fixed \(c_0\in(0,1]\), Theorem~\ref{thm:optimized} and the lower
example give
\begin{equation}
 M_n(c_0;B_n,p_n)\asymp_{c_0,\Lambda}
 \frac{B_n(\log p_n)^{3/2}}{\sqrt n}
 \label{opt:matching-minimax}
\end{equation}
whenever \(p_n\to\infty\), \(B_n\ge2\), the displayed rate tends
to zero, and \(1+\log_+(\sqrt n/B_n)\le\Lambda\log(ep_n)\).
This includes polynomial dimension and growing envelopes subject to
the vanishing-rate condition. The Supplement~\cite[Section~S.4.2,
Theorem~S.4]{Supplement} gives the general consistency argument below.

\begin{theorem}[Uniform consistency at a fixed gap]
\label{opt:consistency}
For fixed \(c_0\in(0,1]\), arbitrary dimension sequences \(p_n\ge1\),
and envelopes \(B_n\ge2\),
\begin{equation}
 M_n(c_0;B_n,p_n)\longrightarrow0
 \quad\Longleftrightarrow\quad
 \frac{B_n\{\log(ep_n)\}^{3/2}}{\sqrt n}\longrightarrow0.
 \label{opt:consistency-iff}
\end{equation}
For each fixed \(B\ge1\),
\[
 M_n^0(c_0;B,p_n)\longrightarrow0
 \quad\Longleftrightarrow\quad \log p_n=o(\sqrt n).
\]
\end{theorem}

For the zero-skewness assertion, independent Rademacher coordinates
locate the critical dimension at
\(\log p\asymp\sqrt n\). If \(\log p_n\sim a\sqrt n\), their
maximum has a limiting discrepancy
\(e^{-\exp(-a^2/3)}-e^{-1}>0\) at a threshold chosen so that the
Gaussian upper-tail count tends to one. The fourth-order tail calculation and the necessity argument for
arbitrary dimension sequences are in the Supplement~\cite[Section~S.3,
Proposition~S.2 and the following proof]{Supplement}. The sharpness of the
threshold says nothing about the error rate below it. In the range
\(n^{1/3}\ll L\ll\sqrt n\), known lower bounds have scale
\(L^2/n\), while the fixed-gap zero-skewness upper bound is larger.
The lattice obstruction must also be retained in any improved
unsmoothed error bound.

For unrestricted covariance the upper bound is optimal up to
logarithmic factors, and where the lower scale \(\delta L^{3/2}\) is
attained, the upper and lower bounds differ by at most the factor
\(\sqrt{D_L}\). It is not known whether this factor is necessary, nor how
the optimal bound depends on a vanishing correlation gap. The companion
note~\cite{Obstruction} shows that natural pointwise strengthenings of
the fourth-order remainder bound, for a two-scale version of the
smoother, are false, so the factor cannot be removed by recovering a
cancellation lost in a pointwise derivative estimate. It also states an
integrated estimate of the remainder that would suffice, and a second
note~\cite{IntegratedRemainder} proves this estimate when the covariance
has rank one.

Exactly repeated or opposite coordinates can be removed without
changing rectangle distance; see the Supplement~\cite[Section~S.5,
Proposition~S.6]{Supplement}. It also gives the rank-one minimax rate
\(C\min\{1,B/\sqrt n\}\) \cite[Section~S.4.3, Proposition~S.4]{Supplement}.
The proof for arbitrary rank does not use these reductions.

Corollary~\ref{cor:subexponential} complements the companion
paper~\cite{Local}, whose sub-exponential bound is sharper in logarithms
when all principal covariance submatrices of order at most four are
uniformly nondegenerate. Without that restriction,
Corollary~\ref{cor:subexponential} still gives the power \(n^{-1/2}\) up
to logarithms. Appendix~\ref{app:subexponential} compares the
assumptions and rates in detail.

% Source: references
\begin{thebibliography}{99}
\phantomsection
\addcontentsline{toc}{section}{References}
\bibitem{AL}
P.~M. Aronow and P. Lopatto (2026).
\emph{Optimal central limit theorem for bounded random variables in high dimensions}.
Manuscript, arXiv:2609.07067v1.
\url{https://arxiv.org/abs/2609.07067v1}.

\bibitem{CCK2013}
V. Chernozhukov, D. Chetverikov, and K. Kato (2013).
\emph{Gaussian approximations and multiplier bootstrap for maxima of
sums of high-dimensional random vectors}.
The Annals of Statistics \textbf{41}(6), 2786--2819.
\href{https://doi.org/10.1214/13-AOS1161}{doi:10.1214/13-AOS1161}.

\bibitem{CCK2017}
V. Chernozhukov, D. Chetverikov, and K. Kato (2017).
\emph{Central limit theorems and bootstrap in high dimensions}.
The Annals of Probability \textbf{45}(4), 2309--2352.
\href{https://doi.org/10.1214/16-AOP1113}{doi:10.1214/16-AOP1113}.

\bibitem{CCKK2022}
V. Chernozhukov, D. Chetverikov, K. Kato, and Y. Koike (2022).
\emph{Improved central limit theorem and bootstrap approximations in
high dimensions}. The Annals of Statistics \textbf{50}(5), 2562--2586.
\href{https://doi.org/10.1214/22-AOS2193}{doi:10.1214/22-AOS2193}.

\bibitem{CCKK2023}
V. Chernozhukov, D. Chetverikov, K. Kato, and Y. Koike (2023).
\emph{High-dimensional data bootstrap}.
Annual Review of Statistics and Its Application \textbf{10}, 427--449.
Working version: arXiv:2205.09691.
\url{https://arxiv.org/abs/2205.09691}.

\bibitem{Local}
V. Chernozhukov, D. Chetverikov, K. Kato, and Y. Koike,
with AI collaborators (2026).
\emph{High-dimensional Gaussian approximation over rectangles under
local nondegeneracy}. Companion working manuscript, 21 October 2026.

\bibitem{Supplement}
V.~Chernozhukov, D.~Chetverikov, K.~Kato, and Y.~Koike,
with AI collaborators (2026).
\emph{Supplementary material: High-dimensional Gaussian approximations
over rectangles without covariance restrictions}.
Working manuscript, 20 October 2026, revision \texttt{bbf1ab99c733}.
\href{https://github.com/VC2015/HDCLTs/blob/bbf1ab99c733a398c181ca8c9f5a243a41f258c6/research-rectangle-clt-singular/Rectangle_CLT_Supplement.pdf}{Supplement}.

\bibitem{CCK}
V. Chernozhukov, D. Chetverikov, and Y. Koike (2023).
\emph{Nearly optimal central limit theorem and bootstrap approximations
in high dimensions}.
The Annals of Applied Probability \textbf{33}(3), 2374--2425.
\href{https://doi.org/10.1214/22-AAP1870}{doi:10.1214/22-AAP1870}.
Working version: arXiv:2012.09513v2, May 12, 2021.
\url{https://arxiv.org/abs/2012.09513v2}.

\bibitem{FKLZ}
X. Fang, Y. Koike, S.-H. Liu, and Y.-K. Zhao (2026).
\emph{High-dimensional central limit theorems by Stein's method in
the degenerate case}. The Annals of Applied Probability, to appear.
arXiv:2305.17365.
\href{https://research.cuhk.edu.hk/en/publications/high-dimensional-central-limit-theorems-by-steins-method-in-the-d/}{Accepted manuscript (2026)}.

\bibitem{GaoZhang}
Z. Gao and W. Zhang (2026).
\emph{Cubic-root Gaussian approximation under unrestricted covariance}.
arXiv:2608.30221v1, 31 August 2026.
\url{https://arxiv.org/abs/2608.30221v1}.

\bibitem{Kane}
D.~M. Kane (2014).
\emph{The average sensitivity of an intersection of half spaces}.
Research in the Mathematical Sciences \textbf{1}, Article~13.
\href{https://doi.org/10.1186/s40687-014-0013-6}{doi:10.1186/s40687-014-0013-6}.
Working version: arXiv:1309.2987v3.
\url{https://arxiv.org/abs/1309.2987v3}.

\bibitem{McFadden1989}
D. McFadden (1989).
\emph{A method of simulated moments for estimation of discrete
response models without numerical integration}.
Econometrica \textbf{57}(5), 995--1026.
\href{https://doi.org/10.2307/1913621}{doi:10.2307/1913621}.

\bibitem{OST}
R. O'Donnell, R.~A. Servedio, and L.-Y. Tan (2019).
Fooling polytopes.
In \emph{Proceedings of the 51st Annual ACM SIGACT Symposium on
Theory of Computing}, 614--625.
\href{https://doi.org/10.1145/3313276.3316321}{doi:10.1145/3313276.3316321}.

\bibitem{Prekopa1973}
A.~Pr\'ekopa (1973).
\emph{On logarithmic concave measures and functions}.
Acta Scientiarum Mathematicarum (Szeged) \textbf{34}, 335--343.
\url{https://rutcor.rutgers.edu/Prekopa/pdf/SCIENT2.pdf}.

\bibitem{IntegratedRemainder}
\emph{Toward an integrated curvature estimate: a proved rank-one case
and an entropy identity for the general problem}.
Companion research note, 24 October 2026.
\href{https://github.com/VC2015/HDCLTs/blob/main/research-rectangle-clt-singular/Rectangle_CLT_Integrated_Remainder.pdf}{\nolinkurl{Rectangle_CLT_Integrated_Remainder.pdf}}.

\bibitem{Obstruction}
\emph{A curvature obstruction for two-scale smoothing}.
Companion research note, 23 October 2026.
\href{https://github.com/VC2015/HDCLTs/blob/main/research-rectangle-clt-singular/Rectangle_CLT_Curvature_Obstruction.pdf}{\nolinkurl{Rectangle_CLT_Curvature_Obstruction.pdf}}.

\end{thebibliography}

\clearpage
\appendix
\numberwithin{equation}{section}
\numberwithin{theorem}{section}
\numberwithin{lemma}{section}
\numberwithin{proposition}{section}
\numberwithin{corollary}{section}
\numberwithin{remark}{section}

\section*{Notation in the appendices}
\addcontentsline{toc}{section}{Notation in the appendices}
The following symbols, most of which Section~\ref{sec:smoothing}
introduces, are used throughout
Appendices~\ref{sec:main-proof}, \ref{appx-sec:smoothing},
and~\ref{appx-sec:support}.
Appendices~\ref{appx-app:concentration} and~E use a few local
symbols, introduced where they occur.

\medskip
\begin{center}
\small
\renewcommand{\arraystretch}{1.15}
\begin{tabular}{@{}p{0.27\textwidth}p{0.68\textwidth}@{}}
\(\delta,\ L,\ c,\ \kappa,\ E_\kappa\) & summand envelope,
\(\log(ep)\), correlation gap, and the parameters
of~\eqref{eq:parameters}\\
\(F_{A,h},\ M_h\) & logistic approximation of \(A\) at scale
\(h\), and \(\sup_A|\E F_{A,h}(W)-\E F_{A,h}(Z)|\),
see~\eqref{eq:selector} and~\eqref{eq:smoothed-discrepancy}\\
\(P_A(\lambda),\ \tau_h\) & \(\E F_{A,h}(U_\lambda)\) along
the path~\eqref{eq:path}, and the time \(Lh^2\) at which the early and
late bounds meet\\
\(\nu_\alpha,\ a_\alpha,\ m\) & face normal \(\pm e_j\) and threshold of
face \(\alpha\); there are \(m\le2p\) finite faces\\
\(A^t\) & the rectangle \(A\) with every finite endpoint moved outward
by \(t\), or inward if \(t<0\)\\
\(F,\ \Psi,\ w_\alpha,\ T,\ N\) & logistic approximation, its potential
\(h\log(1+T)\), the face weights, their sum
\(T=\sum_\alpha w_\alpha\), and the denominator \(N=1+T\) of \(F\)\\
\(\pi_\alpha,\ \theta,\ \bar\nu,\ \chi\) & normalized weights, total
mass \(1-F\), weighted face normal \(\nabla\Psi\), and curvature
\(h\tr D^2\Psi\)\\
\(\omega_0,\ \omega_\alpha\) & independent unit exponentials in the
mixture representation of \(F\)\\
\(\Gamma,\ r,\ u_\alpha,\ \norm\cdot_U\) & support factor
\(\Sigma=\Gamma\Gamma^\top\), its rank, face normals
\(\Gamma^\top\nu_\alpha\), and the face norm
\(\max_\alpha|u_\alpha\cdot x|\)\\
\(\zeta_i^0;\ \zeta_i,\ \widetilde\zeta_i\) & summands in support coordinates; two
independent copies of them\\
\(\gamma,\ \varpi,\ \mathbb B_{i,\lambda}\) & amplitude \(1+\lambda\), selection
probability \((1-\lambda)/(1+\lambda)^2\), and the selection variables\\
\(Y_i,\ V_\lambda,\ U_\lambda,\ H_i\) & selected summands
\(\gamma \mathbb B_{i,\lambda}\zeta_i^0\), their sum, the path, and the path with
summand \(i\) deleted\\
\(V_0,\ \sigma,\ G,\ U\) & concentrated base, Gaussian scale, standard
Gaussian vector on \(\R^r\), and \(U=V_0+\sigma G\); the coordinate
standard deviations \(\sigma_j\) of Corollary~\ref{cor:subexponential}
carry a subscript\\
\(\Phi_r,\ \varphi_r,\ H_2,\ \surf\) & standard Gaussian measure and
density on \(\R^r\), \(H_2(g)=\norm g^2-r\), and the Gaussian surface
integral\\
\(\mathsf A,\ \mathsf A_b,\ d_0\) & large absolute constants in
\(h\ge\mathsf A\delta\), \(h\ge\mathsf A_b\delta\), and
\(\lambda\ge d_0\delta^2L\)\\
\(C,\ C_q;\ c_1,\ c_2,\ c_*\) & absolute constants, large and small;
they may change between occurrences, and a subscript records any
dependence. The letter \(c\) alone is the correlation gap.\\
\end{tabular}
\end{center}

% Source: complete_integration_proofs
\section{\texorpdfstring{Proofs of Proposition~\ref{prop:moving-cutoff} and the theorems}{Proofs of Proposition 2 and the theorems}}\label{sec:main-proof}
This appendix proves Proposition~\ref{prop:moving-cutoff} and
Theorems~\ref{thm:unrestricted}--\ref{thm:zero-skew} along the lines of
Subsection~\ref{sec:rate}. It uses only Lemma~\ref{lem:endpoint} and
Proposition~\ref{an:prop:derivative}, which
Appendices~\ref{appx-sec:smoothing} and~\ref{appx-sec:support} prove.
Recall that \(P_A(\lambda)=\E F_{A,h}(U_\lambda)\) along the
path~\eqref{eq:path}, so that \(P_A(0)=\E F_{A,h}(W)\) and
\(P_A(1)=\E F_{A,h}(Z)\), and that
\(M_h=\sup_A|P_A(1)-P_A(0)|\) by~\eqref{eq:smoothed-discrepancy}.

\begin{proof}[Proof of~\eqref{eq:smooth-moving} in
Proposition~\ref{prop:moving-cutoff}]
Increase \(\Ac\) so that \eqref{an:early} and \eqref{an:late} apply and
\(\Ac^2\ge d_0\). Put \(\tau_h=Lh^2\), the time at which their
fourth-order terms agree. Since \(h\ge\Ac\delta\), we have
\(\tau_h\ge d_0\delta^2L\). If \(\tau_h<1\), the later estimate
is therefore available throughout \([\tau_h,1]\). Integrating the
early estimate gives
\begin{align*}
 |P_A(\tau_h)-P_A(0)|
 &\le C\sqrt L\left(
 \frac{\delta\tau_h^2}{h^2}
 +\frac{\delta^2\tau_h}{h^3}\right)\\
 &=C\delta L^{3/2}
 \left(Lh^2+\frac\delta h\right)
 \le C\delta L^{3/2}.
\end{align*}
The later estimate gives
\[
 |P_A(1)-P_A(\tau_h)|
 \le CL^{3/2}\left[
 \delta+\frac{\delta^2}{h}\log\frac1{\tau_h}\right].
\]
The constant cubic contribution integrates to \(C\delta L^{3/2}\);
the fourth-order factor \(1/\lambda\) gives the logarithm. Combining
the two time ranges and taking the supremum over thresholds proves
\eqref{eq:smooth-moving} when \(\tau_h<1\).

If \(Lh^2\ge1\), integrate the early estimate over the whole path:
\[
 |P_A(1)-P_A(0)|
 \le C\sqrt L\left(\frac\delta{h^2}
 +\frac{\delta^2}{h^3}\right)
 =C\delta L^{3/2}\frac{1+\delta/h}{Lh^2}
 \le C\delta L^{3/2}.
\]
The logarithmic term in \eqref{eq:smooth-moving} is zero in this case.
Continuity at the interpolation endpoints permits integration up to
\(0\) and \(1\). All bounds are uniform in the thresholds.
\end{proof}

\begin{proof}[Proof of Theorem~\ref{thm:optimized}]
Combining Proposition~\ref{prop:moving-cutoff} with
\eqref{eq:endpoint-min} yields, for every \(h\ge\Ac\delta\),
\begin{equation}
 \rh(W,Z)\le C\left[
 \delta L^{3/2}
 +\frac{\delta^2}{h}L^{3/2}\log_+\frac1{Lh^2}
 +h\sqrt L\,\kappa\right].
 \label{eq:variational}
\end{equation}
The first term is independent of \(h\). To balance the fourth-order
comparison error against the endpoint error, choose \(h\) as
in~\eqref{eq:optimal-scale}. Because \(\kappa\le L\) and \(E_\kappa\ge1\), this scale is at least
\(\Ac\delta\). Moreover,
\[
 \log_+\frac1{Lh^2}
 =2\log_+\frac{\sqrt\kappa}{\Ac\delta L\sqrt{E_\kappa}}
 \le2\log_+\frac{\sqrt\kappa}{\delta L}
 \le2E_\kappa.
\]
The endpoint term in \eqref{eq:variational} equals
\(\Ac\delta L\sqrt{\kappa E_\kappa}\), while its logarithmic term is at
most \((2/\Ac)\delta L\sqrt{\kappa E_\kappa}\).
This proves \eqref{eq:optimized} after taking the minimum with one.
The version with \(D_L\) follows from \(E_\kappa\le D_L\).
\end{proof}

\begin{proof}[Proof of Theorem~\ref{thm:unrestricted}]
Since \(\kappa\le L\), \(E_\kappa\le D_L\), and \(D_L\ge1\),
\[
 L^{3/2}+L\sqrt{\kappa E_\kappa}
 \le 2L^{3/2}\sqrt{D_L}.
\]
Apply Theorem~\ref{thm:optimized}.
\end{proof}

% Source: zero_skewness_proof
\subsection{\texorpdfstring{Third-moment cancellation: the second part of Proposition~\ref{prop:moving-cutoff} and Theorem~\ref{thm:zero-skew}}{Third-moment cancellation: the second part of Proposition 2 and Theorem 3}}\label{zs:section}
This subsection proves the second part of
Proposition~\ref{prop:moving-cutoff} and Theorem~\ref{thm:zero-skew}.

\begin{proof}[Proof of~\eqref{zs:smooth} in
Proposition~\ref{prop:moving-cutoff}]
Use \eqref{zs:identity} to cancel the common-base cubic term before
estimating the derivative. The separate fourth-order bounds
\eqref{eq:remainder-early} and \eqref{an:remainder} then give
\begin{align}
 |P_A'(\lambda)|
 &\le \frac{C\delta^2\sqrt L}{h^3},
 &&0<\lambda<1,\label{zs:early}\\
 |P_A'(\lambda)|
 &\le \frac{C\delta^2L^{3/2}}{\lambda h},
 &&d_0\delta^2L\le\lambda<1.\label{zs:late}
\end{align}
Increase \(\Ac\) so that \(\Ac^2\ge d_0\), and put \(\tau_h=Lh^2\).
If \(\tau_h<1\), then \(\tau_h\ge d_0\delta^2L\). Integrating
\eqref{zs:early} up to \(\tau_h\), and \eqref{zs:late} after it,
yields
\begin{align*}
 |P_A(1)-P_A(0)|
 &\le \frac{C\delta^2\sqrt L}{h^3}\tau_h
       +\frac{C\delta^2L^{3/2}}h\log\frac1{\tau_h}\\
 &=\frac{C\delta^2L^{3/2}}h
       \left[1+\log\frac1{Lh^2}\right].
\end{align*}
If \(Lh^2\ge1\), integrate \eqref{zs:early} on the whole path and
use \(\sqrt L/h^3\le L^{3/2}/h\). Endpoint continuity, already
established after Proposition~\ref{an:prop:derivative}, justifies
integration to \(0\) and \(1\). Taking the supremum over thresholds
proves \eqref{zs:smooth}.
\end{proof}

\begin{proof}[Proof of Theorem~\ref{thm:zero-skew}]
Combine the second part of Proposition~\ref{prop:moving-cutoff} with the endpoint
bound \eqref{eq:endpoint-min} to obtain, for every \(h\ge\Ac\delta\),
\begin{equation}
 \rh(W,Z)\le C\left\{
 \frac{\delta^2L^{3/2}}h
 \left[1+\log_+\frac1{Lh^2}\right]+h\sqrt L\,\kappa\right\}.
 \label{zs:variational}
\end{equation}
To balance the comparison error against the endpoint error, choose
\(h\) by~\eqref{eq:optimal-scale} again; it is admissible because \(\kappa\le L\) and \(E_\kappa\ge1\).
As before,
\[
 \log_+\frac1{Lh^2}
 =2\log_+\frac{\sqrt\kappa}{\Ac\delta L\sqrt{E_\kappa}}
 \le2E_\kappa.
\]
The endpoint term in \eqref{zs:variational} equals
\(\Ac\delta L\sqrt{\kappa E_\kappa}\); its comparison term is at most
\[
 \frac{\delta L\sqrt\kappa}{\Ac\sqrt{E_\kappa}}(1+2E_\kappa)
 \le\frac3{\Ac}\delta L\sqrt{\kappa E_\kappa}.
\]
Taking the minimum with one proves \eqref{eq:zero-skew}.

\end{proof}

% Source: smoothing_proofs
\section{\texorpdfstring{The logistic approximation: proof of Lemma~\ref{lem:endpoint}}{The logistic approximation: proof of Lemma 2}}\label{appx-sec:smoothing}

This appendix develops the properties of the logistic
approximation \(F_{A,h}\) of Subsection~\ref{sec:outline-smoothing} that
the proofs use. We first record how Lemma~\ref{lem:concentration}
applies to systems of halfspaces, then derive the mixture and derivative
formulas for \(F_{A,h}\), and finally bound the error that it introduces
at \(W\) and \(Z\), which proves Lemma~\ref{lem:endpoint}. The
derivative formulas and the averaged envelope of
Lemma~\ref{appx-lem:envelope} are used again in
Appendix~\ref{appx-sec:support}.

\subsection{\texorpdfstring{Anti-concentration for systems of halfspaces}{Anti-concentration for systems of halfspaces}}

For a fixed system of halfspaces, the notation
\(A^{t}\) of Subsection~\ref{sec:outline-smoothing} increases all its
thresholds by \(t\). We also apply Lemma~\ref{lem:concentration},
which Appendix~\ref{appx-app:concentration} proves, to the
vector of face projections
\((\nu_\alpha\cdot x)_{\alpha\le m}\). Its coordinates retain
the same variance and summand bound. The intersection
\(\{\nu_\alpha\cdot x\le a_\alpha\}_{\alpha\le m}\) is the inverse
image of the lower orthant \(\prod_{\alpha\le m}(-\infty,a_\alpha]\)
in \(\R^m\); comparing thresholds increased and decreased by \(t\)
gives the boundary layer in the concentration estimate. The orthant is nonempty even when
its inverse image is empty, so systems of inequalities with no
common solution are covered. Since \(m\le2p\), we have \(\log(em)\le CL\).

\subsection{\texorpdfstring{Mixture and derivative formulas}{Mixture and derivative formulas}}\label{appx-sec:mixture}

We keep the notation of Subsection~\ref{sec:outline-smoothing} for
the faces \(\alpha=1,\ldots,m\) of \(A\), with face normals
\(\nu_\alpha\) and thresholds \(a_\alpha\), for the face weights
\(w_\alpha\), the logistic approximation \(F=F_{A,h}\), the face
probabilities \(\pi_\alpha\), and \(\theta=1-F\).
Recall that \(T=\sum_\alpha w_\alpha\) and
\(\Psi=\Psi_{A,h}=h\log(1+T)\). In addition, put
\[
 \pi_0=F=\frac1{1+T},\qquad
 \bar\nu=\sum_\alpha\pi_\alpha\nu_\alpha,
 \qquad \chi=h\,\tr D^2\Psi.
\]

Here \(\pi_0\) is the probability of the baseline category, the
vector \(\bar\nu=\nabla\Psi\) is the weighted face normal, and
the scalar \(\chi\) measures the curvature used in the gap-sensitive
endpoint estimate below. We use these formulas for arbitrary
coordinate thresholds, including those with an empty intersection,
because the shifts below need not preserve the ordering of the
endpoints. In such cases, \(F_{A,h}\) refers to the specified
inequalities, and its denominator is positive. If \(m=0\), the logistic
approximation is identically one and all comparisons are trivial; the
analytic arguments below take \(m\ge1\).

The mixture representation~\eqref{eq:mixture} expresses \(F\) as
an average of indicators of the systems \(A_{\omega}\), whose
thresholds are shifted at random by
\(h\log(\omega_\alpha/\omega_0)\), where
\(\omega_0,(\omega_\alpha)_\alpha\) are independent unit
exponentials. To verify it, observe that, conditional on
\(\omega_0\), the probability that \(x\in A_{\omega}\) is
\(e^{-\omega_0T(x)}\), and its expectation is \((1+T(x))^{-1}\).
Equivalently, the variables \(-\log\omega_\beta\) are independent standard
Gumbel variables, so
\[
 F_{A,h}(x)=\Pp\!\left\{-h\log\omega_0\ge
 \max_\alpha(\nu_\alpha\cdot x-a_\alpha-h\log\omega_\alpha)\right\}.
\]
Here the baseline category has deterministic score zero, and the face
displacements \(h\log(\omega_\alpha/\omega_0)\) have logistic marginals and are
dependent through \(\omega_0\); independent logistic displacements would give
a different approximation. McFadden~\cite[p.~1001]{McFadden1989} uses
the same random-utility representation, it is stated for softmax weights in~\cite[Lemma~14]{GaoZhang}. 

The general derivative formula described in
Subsection~\ref{sec:outline-smoothing} reads, for vectors
\(z_1,\ldots,z_q\),
\begin{equation}
 \begin{split}
 D^qF(x)[z_1,\ldots,z_q]
 ={}&h^{-q}\sum_{\mathcal P\in\Pi_q}(-1)^{|\mathcal P|}|\mathcal P|!
 \sum_{\alpha_1,\ldots,\alpha_{|\mathcal P|}}
 F(x)\prod_{b=1}^{|\mathcal P|}\pi_{\alpha_b}(x)\\
 &\hspace{3em}\times
 \prod_{b=1}^{|\mathcal P|}\prod_{a\in J_b}
 (\nu_{\alpha_b}\cdot z_a),
 \end{split}
 \label{appx-eq:partition-derivative}
\end{equation}
where \(\Pi_q\) is the set of partitions \(\mathcal P\) of \(\{1,\ldots,q\}\),
and \(J_1,\ldots,J_{|\mathcal P|}\) are the blocks in any fixed ordering.
Ordered partitions can be encoded by a permutation and cuts between
successive elements, so
\(\sum_{\mathcal P\in\Pi_q}|\mathcal P|!\le q!2^{q-1}\).
In particular,
\begin{equation}
 |D^qF(x)[z_1,\ldots,z_q]|
 \le C^q q!h^{-q}\prod_{a=1}^q
 \max_\alpha|\nu_\alpha\cdot z_a|.
 \label{appx-eq:all-orders}
\end{equation}
The softmax cumulant formula~\cite[Lemma~13]{GaoZhang} derives
\eqref{appx-eq:partition-derivative} and~\eqref{appx-eq:all-orders}
from the joint cumulants of a categorical variable.

The uniform bound \eqref{appx-eq:all-orders} does not yet record that the
derivatives are small away from the boundary. For the ordered
coordinate tuples \(I\) and the derivatives \(\partial_I\) of
Subsection~\ref{sec:outline-smoothing}, the same positive majorant,
using \(\norm{\nu_\alpha}_1=1\), gives
\begin{equation}
 \sum_I\sup_{\norm z_\infty\le\varrho}|\partial_IF(x+z)|
 \le C_qh^{-q}F(x)(1-F(x)),\qquad 0\le\varrho\le h.
 \label{appx-eq:local-envelope}
\end{equation}
To see the local stability, observe that each \(w_\alpha\)
changes by a factor in
\([e^{-\varrho/h},e^{\varrho/h}]\), and \(1+T\) obeys the same comparison.
After summing over coordinate tuples and face indices, the terms with
\(k\) blocks are bounded by a constant times
\(h^{-q}T^k/(1+T)^{k+1}\le h^{-q}F(1-F)\).

In the next lemma, \(U\) has an independent component in the sense
that \(U=S+R\), where \(S\) is independent of \(R\) and \(S\) is either a
sum of independent centered vectors with summand bound \(2\delta\) and
equal coordinate variances at least \(1/2\), or a centered Gaussian vector
with equal coordinate variances at least \(1/2\).

\begin{lemma}[Averaging the derivative envelope]\label{appx-lem:envelope}
Suppose a random vector \(U\) has an independent component that is
either a sum of independent centered vectors, with summand bound
\(2\delta\) and equal coordinate variances at least \(1/2\), or a
centered Gaussian vector with equal coordinate variances at least
\(1/2\).
Uniformly over thresholds, for \(h\ge4\delta\),
\begin{equation}
 \E F(U)(1-F(U))\le Ch\sqrt L.
 \label{appx-eq:small-envelope}
\end{equation}
Consequently, for fixed \(q\) and \(h\ge\max\{4\delta,\varrho\}\),
\begin{equation}
 \E\sum_I\sup_{\norm z_\infty\le\varrho}|\partial_IF(U+z)|
 \le C_q h^{1-q}\sqrt L.
 \label{appx-eq:averaged-derivatives}
\end{equation}
\end{lemma}

\begin{proof}
Shift every face threshold by the same amount \(s\), and write
\(F_s=(1+e^{-s/h}T)^{-1}\) for the resulting logistic approximation. Since
\(\partial_sF_s=F_s(1-F_s)/h\), and
\(F_s(1-F_s)\ge e^{-1}F(1-F)\) for \(|s|\le h\),
\[
 F(1-F)\le\frac e2(F_h-F_{-h}).
\]
In the mixture representation~\eqref{eq:mixture},
\(F_h-F_{-h}\) is the average of
\(\ind_{A_{\omega}^h}-\ind_{A_{\omega}^{-h}}\), with shifts interpreted in the
face projections \(\nu_\alpha\cdot x\). Conditional on the exponentials and the
independent shift, Lemma~\ref{lem:concentration} bounds the expectation
by \(C(h+\delta)\sqrt L\), since the base has equal coordinate
variances
at least \(1/2\) and summand bound at most \(2\delta\), or is
Gaussian. Since \(h\ge4\delta\), this proves
\eqref{appx-eq:small-envelope}. Combining it with
\eqref{appx-eq:local-envelope} proves \eqref{appx-eq:averaged-derivatives}.
\end{proof}

\subsection{The error at the endpoints}

The mixture representation gives an endpoint bound for every
covariance matrix, and a second argument, which uses the correlation
gap, can be sharper. We choose the better of the two after fixing the
smoothing scale.

\begin{lemma}[Unrestricted endpoint]\label{appx-lem:endpoint-free}
Under \eqref{eq:setup}, for \(h\ge\delta\),
\begin{equation}
 \E|F_{A,h}(W)-\ind_A(W)|
 +\E|F_{A,h}(Z)-\ind_A(Z)|
 \le ChL^{3/2}.
 \label{appx-eq:endpoint-free}
\end{equation}
\end{lemma}

\begin{proof}
Recall the largest threshold change
\(R_h=h\max_\alpha|\log(\omega_\alpha/\omega_0)|\)
from Subsection~\ref{sec:outline-smoothing}.
Then \(A^{-R_h}\subseteq A_{\omega}\subseteq A^{R_h}\), and
\(\E R_h\le Ch\log(em)\le ChL\).
The last assertion follows by integrating the exponential tails of
\(-\log\omega_\alpha\), and the faster upper tails of \(\log\omega_\alpha\).
For fixed exponentials, the indicators can disagree only in
\(A^{R_h}\setminus A^{-R_h}\). Jensen's inequality moves the absolute
value inside the exponential average. Conditional on the exponentials,
Lemma~\ref{lem:concentration} bounds the errors at \(W\) and \(Z\)
by \(C(R_h+\delta)\sqrt L\) and \(CR_h\sqrt L\), respectively.
Averaging the radius gives \(C(hL+\delta)\sqrt L\) and
\(ChL\sqrt L\).
Use \(h\ge\delta\) and \(L\ge1\).
\end{proof}

\begin{lemma}[Curvature endpoint]\label{appx-lem:endpoint-gap}
Under \eqref{eq:setup}, if \(c>0\) and \(h\ge\delta\), then
\begin{equation}
 \E|F_{A,h}(W)-\ind_A(W)|
 +\E|F_{A,h}(Z)-\ind_A(Z)|
 \le \frac{Ch\sqrt L}{c}.
 \label{appx-eq:endpoint-gap}
\end{equation}
\end{lemma}

\begin{proof}
We first bound the error of \(F\) pointwise by the curvature
\(\chi\) of \(\Psi\), and then bound the expected curvature by means of
the correlation gap. Direct differentiation gives
\begin{equation}
 D^2\Psi=\frac1h\left(
 \sum_\alpha\pi_\alpha\nu_\alpha\nu_\alpha^\top-\bar\nu\bar\nu^\top\right),
 \qquad \norm{\bar\nu}_1\le\theta,\qquad
 \chi=\theta-\norm{\bar\nu}_2^2.
 \label{appx-eq:hessian}
\end{equation}
Distinct face normals \(\pm e_j\) have nonpositive inner products.
Thus, with \(\pi_* =\max_\alpha\pi_\alpha\),
\(\chi\ge\theta(1-\theta)\) and
\(\chi\ge\theta(1-\pi_*)\).
On \(A\), every \(w_\alpha\le1\), so \(\pi_*\le1/2\).
Off \(A\) we have \(T>1\), so \(\theta>1/2\). On \(A\), the second
lower bound gives \(\chi\ge\theta/2\); off \(A\), the first gives
\(\chi\ge F/2\). The indicator error is \(\theta\) and \(F\)
in the two cases, respectively, so
\begin{equation}
 |F-\ind_A|\le2\chi.
 \label{appx-eq:curvature-error}
\end{equation}
To bring the covariance into this estimate, observe that
\[
 \sum_{j\ne k}|\partial_{jk}\Psi|
 =h^{-1}(\norm{\bar\nu}_1^2-\norm{\bar\nu}_2^2)
 \le\tr D^2\Psi,
\]
and since \(\Sigma_{jj}=1\) and \(|\Sigma_{jk}|\le1-c\) for
\(j\ne k\), the off-diagonal contribution can remove at most
\((1-c)\tr D^2\Psi\). Therefore
\begin{equation}
 \tr(\Sigma D^2\Psi)\ge c\,\tr D^2\Psi.
 \label{appx-eq:gap-contraction}
\end{equation}

To compare the Hessian at nearby points, include the index \(0\), with
\(\nu_0=0\) and \(\pi_0=F\); the positive representation
\[
 D^2\Psi=\frac1{2h}\sum_{\alpha,\beta\in\{0,1,\ldots,m\}}
 \pi_\alpha\pi_\beta
 (\nu_\alpha-\nu_\beta)(\nu_\alpha-\nu_\beta)^\top
\]
implies
\(e^{-4\varrho/h}D^2\Psi(x)\preceq D^2\Psi(x+z)
\preceq e^{4\varrho/h}D^2\Psi(x)\) when \(\norm z_\infty\le\varrho\);
here \(M\preceq N\) means that \(N-M\) is positive semidefinite.
Set \(W_i=W-\xi_i\). Centering and independence let us subtract
\(\nabla\Psi(W_i)\) before integrating the Hessian along the segment
from \(W_i\) to \(W\). First compare this Hessian with
\(D^2\Psi(W_i)\). Independence then replaces
\(\xi_i\xi_i^\top\) by \(\E\xi_i\xi_i^\top\). A second comparison
replaces \(D^2\Psi(W_i)\) by \(D^2\Psi(W)\), and summing the covariances of the
summands gives \(\Sigma\). Each lower comparison contributes a factor
\(e^{-4\delta/h}\), so
\[
 \begin{split}
 \E[W\cdot\nabla\Psi(W)]
 &=\sum_i\E\int_0^1
 \xi_i^\top D^2\Psi(W_i+s\xi_i)\xi_i\dd s\\
 &\ge e^{-8\delta/h}\E\tr(\Sigma D^2\Psi(W)).
 \end{split}
\]
Since \(\norm{\nabla\Psi}_1\le1\), the left side is at most
\(\E\norm W_\infty\), which Bernstein's inequality bounds by
\(C(\sqrt L+\delta L)\). When \(\delta\sqrt L\le1\), this is
\(C\sqrt L\), and \(h\ge\delta\) bounds \(e^{8\delta/h}\)
absolutely. Combining with \eqref{appx-eq:gap-contraction} gives
\(\E\chi(W)\le Ch\sqrt L/c\).
If \(\delta\sqrt L>1\), the same conclusion follows from
\(\chi\le1\) and \(h\ge\delta\).
For \(Z\), Gaussian integration by parts gives
\[
 \E[Z\cdot\nabla\Psi(Z)]
 =\E\tr(\Sigma D^2\Psi(Z)),\qquad
 \E\norm Z_\infty\le C\sqrt L.
\]
This identity holds also for singular \(\Sigma\), as one sees by
restricting to its support. Use \eqref{appx-eq:curvature-error} to finish.
\end{proof}

\begin{proof}[Proof of Lemma~\ref{lem:endpoint}]
Lemma~\ref{appx-lem:endpoint-free} bounds the endpoint error by
\(ChL^{3/2}=Ch\sqrt L\cdot L\), and, when \(c>0\),
Lemma~\ref{appx-lem:endpoint-gap} bounds it by \(Ch\sqrt L\cdot c^{-1}\).
The smaller of the two is \(Ch\sqrt L\,\kappa\), which
is~\eqref{eq:endpoint-error}. Since
\(|\Pp(W\in A)-\Pp(Z\in A)|\) is at most
\(|\E F_{A,h}(W)-\E F_{A,h}(Z)|\) plus the two endpoint errors,
\eqref{eq:endpoint-min} follows.
\end{proof}

% Replace the complete Appendix C (including its old outer revblock).
% This file uses the existing manuscript preamble and bibliography.
% Prepared against Git blob 2076dea4e857bbd02277b4a6f5f92135ace8a065.
% Insert together with Rectangle_CLT_Section_2-reworking-v2.tex.
\section{Comparison along the path: proof of Proposition~\ref{an:prop:derivative}}
\label{appx-sec:support}

The early and late bounds estimate the same third- and fourth-order
terms. Their common ingredient is the covariance identity for the
summands: it controls a pair of directions in each derivative without
introducing the covariance rank. Near the bounded endpoint, we use
the local derivative envelope. Once enough Gaussian variance is
present, we move that pair of derivatives onto the Gaussian density.

We first choose coordinates on the covariance support and derive the
exact interpolation identity. Subsections~\ref{appx-an:sec:early}
and~\ref{appx-an:sec:restoration} then prove the two bounds, using
the matrix estimate of Subsection~\ref{appx-an:sec:matrix} for the
late bound. Subsection~\ref{appx-an:sec:surface} supplies the surface
estimate needed in that matrix argument. We assume that at least one
face is present; otherwise \(F\equiv1\) and the comparison is trivial.

\subsection{Coordinates on the covariance support}
\label{appx-an:sec:coords}

Factor \(\Sigma=\Gamma\Gamma^\top\), where \(\Gamma\) has full column
rank \(r=\operatorname{rank}\Sigma\). Every summand lies almost surely
in its range: for each null direction \(z\) of \(\Sigma\),
\(\sum_i\E(z\cdot\xi_i)^2=z^\top\Sigma z=0\).
We may therefore write \(\xi_i=\Gamma\zeta_i^0\), and injectivity of
\(\Gamma\) gives \(\sum_i\E\zeta_i^0(\zeta_i^0)^\top=I_r\).
For each face normal \(\nu_j\), put \(u_j=\Gamma^\top\nu_j\).
Unit marginal variances give \(\norm{u_j}=1\). The two properties
of these coordinates that we will use are
\begin{equation}
 \sum_i\E\zeta_i^0(\zeta_i^0)^\top=I_r,
 \qquad \norm{\zeta_i^0}_U\le\delta,
 \qquad \norm x_U:=\max_{j\le m}|u_j\cdot x|.
 \label{appx-eq:projected-bound}
\end{equation}
The second follows from \(u_j\cdot\zeta_i^0=\nu_j\cdot\xi_i\).
Thus small positive eigenvalues of \(\Sigma\) do not affect the
quantities used in the proof. This is the support-coordinate reduction
of~\cite[Section~2.3]{GaoZhang}.

We keep the notation \(F(x)=F_{A,h}(\Gamma x)\) and write
\[
 t_j(x)=(u_j\cdot x-a_j)/h,\quad N(x)=1+\sum_j e^{t_j(x)},\quad
 F=N^{-1},\quad \pi_j=e^{t_j}/N,\quad\theta=1-F.
\]
The derivative formulas of Appendix~\ref{appx-sec:smoothing} apply
with \(u_j\) in place of \(\nu_j\).

The covariance identity controls products of two face projections.
We will also use its conditional version. If a random array \((Z_i)_i\)
and a \(\sigma\)-field \(\mathcal H\) satisfy
\(\sum_i\E[Z_iZ_i^\top\mid\mathcal H]\preceq I_r\), then, for
any unit vectors \(u,v\), conditional Cauchy--Schwarz gives
\begin{equation}
 \sum_i\E[|(u\cdot Z_i)(v\cdot Z_i)|\mid\mathcal H]
 \le\left\{\sum_i\E[(u\cdot Z_i)^2\mid\mathcal H]
           \sum_i\E[(v\cdot Z_i)^2\mid\mathcal H]\right\}^{1/2}
 \le1.
 \label{appx-an:pair-contraction}
\end{equation}
Bounding the other directions by \(\delta\) will therefore cost
\(\delta\) for a third derivative and \(\delta^2\) for a fourth
derivative. The same covariance bound will later control coefficient
matrices in operator norm.

We also record the concentration property that supplies the small
boundary factor. For \(a\in\R^m\), define
\[
 Q(a)=\{x:u_j\cdot x\le a_j,\ j\le m\},\qquad
 Q(a)^t=Q(a+t\mathbf1).
\]
We call \(V_0\) a \emph{concentrated base} if, for an absolute
constant \(C_0\),
\begin{equation}
 \Pp\{V_0\in Q(a)^t\setminus Q(a)^{-t}\}
 \le C_0(t+\delta)\sqrt L,
 \qquad a\in\R^m,\quad t>0.
 \label{appx-an:concentration}
\end{equation}
Empty intersections are allowed. The property passes to subsystems
by sending unused thresholds to infinity, and to independent shifts
by conditioning. Lemma~\ref{lem:concentration}, applied to the face
projections, establishes it whenever an independent component is
either a bounded sum with summand face bound \(C\delta\), or a
Gaussian vector, and its equal face variances are bounded below by
a positive absolute constant.

For a concentrated base, the envelope argument of
Lemma~\ref{appx-lem:envelope} gives, when \(h\ge\delta\),
\begin{equation}
 \sup_z\E F(V_0+z)\theta(V_0+z)\le Ch\sqrt L.
 \label{appx-an:weight-concentration}
\end{equation}
Indeed, shifting all thresholds by \(\pm h\) gives
\(F(1-F)\le(e/2)(F_h-F_{-h})\). The mixture representation writes
the difference as an average of boundary layers, each controlled by
\eqref{appx-an:concentration}. The same argument for one face bounds
\(e^{t_j}/(1+e^{t_j})^2\). We use these facts after conditioning on
an independent Gaussian shift; no conditional concentration of
\(V_0\) itself is required.

\subsection{The path and its derivative}
\label{appx-an:sec:path}

In support coordinates the path is
\begin{equation}
 Y_i=\gamma\mathbb B_{i,\lambda}\zeta_i^0,\quad
 V_\lambda=\sum_iY_i,\quad U_\lambda=V_\lambda+\sqrt\lambda\,G,
 \quad \gamma=1+\lambda,\quad
 \varpi=\frac{1-\lambda}{(1+\lambda)^2},
 \label{appx-an:path}
\end{equation}
where \(\mathbb B_{i,\lambda}\sim\operatorname{Bernoulli}(\varpi)\)
and \(G\sim N(0,I_r)\) are independent of the summands and of one
another. Only the law at each \(\lambda\) is used. Since
\(\varpi\gamma^2=1-\lambda\), the covariance of \(U_\lambda\)
is \(I_r\), and \(P_A(\lambda)=\E F(U_\lambda)\).
The additional calibration \(\varpi\gamma^3=1-\lambda^2\)
makes the cubic coefficient stationary at \(\lambda=0\).

Let \(H_i=U_\lambda-Y_i\). Let \(\zeta_i,\widetilde\zeta_i\)
be independent fresh copies of \(\zeta_i^0\), independent of the
entire path. We abbreviate repeated arguments by
\(D^qF(x)[z^q]\). To see the cancellations before expanding the
remainder, define the signed density
\[
 \nu_\lambda(s)=\frac{4-2(3-\lambda)s}{\gamma},\qquad 0\le s\le1.
\]
Then
\begin{equation}
 P_A'(\lambda)=\frac12\sum_i\E\left[
 D^2F(U_\lambda)[\zeta_i^2]
 -\int_0^1\nu_\lambda(s)
       D^2F(H_i+s\gamma\zeta_i)[\zeta_i^2]\dd s\right].
 \label{appx-an:signed-hessian}
\end{equation}
Its two relevant moments are
\begin{equation}
 \int_0^1\nu_\lambda(s)\dd s=1,\qquad
 \int_0^1s\nu_\lambda(s)\dd s=\frac{2\lambda}{3\gamma}.
 \label{appx-an:signed-moments}
\end{equation}
The mass one cancels the constant Hessian term; the first moment
produces the coefficient \(-\lambda/3\) of the cubic term.
The density may be negative, and should not be interpreted as a
probability density. Its total variation is at most \(5/3\).

To prove~\eqref{appx-an:signed-hessian}, differentiate one Bernoulli
law at a time, then its amplitude and the Gaussian variance. This gives
\begin{align*}
 P_A'(\lambda)
 ={}&\frac12\sum_i\E D^2F(U_\lambda)[\zeta_i^2]\\
 &+\sum_i\left[\varpi'\E\{F(H_i+\gamma\zeta_i)-F(H_i)\}
             +\varpi\E DF(H_i+\gamma\zeta_i)[\zeta_i]\right].
\end{align*}
Here the first term is the Gaussian heat derivative, using
\(\sum_i\E\zeta_i\zeta_i^\top=I_r\). At a fixed base \(x\),
centering and the fundamental theorem of calculus give
\begin{align*}
 \E\{F(x+\gamma\zeta_i)-F(x)\}
 &=\gamma^2\int_0^1(1-s)
       \E D^2F(x+s\gamma\zeta_i)[\zeta_i^2]\dd s,\\
 \E DF(x+\gamma\zeta_i)[\zeta_i]
 &=\gamma\int_0^1\E D^2F(x+s\gamma\zeta_i)[\zeta_i^2]\dd s.
\end{align*}
Since \(\varpi'=-(3-\lambda)/\gamma^3\), their combined coefficient
is \(\varpi'\gamma^2(1-s)+\varpi\gamma=-\nu_\lambda(s)/2\).
This proves~\eqref{appx-an:signed-hessian}; integrating the affine
density gives~\eqref{appx-an:signed-moments}. All differentiations
are justified for fixed \(h>0\) by bounded summands and bounded
derivatives of \(F\).

We now put the cubic term at the full path, where aggregate third
moments can cancel. The three corrections below respectively record
the Hessian expansion along a fresh summand, the difference between
full and deleted bases for the quadratic term, and the restoration
of the cubic term to the full path. The kernel agrees with
\cite[Lemma~4.1]{AL}, and the resulting full-path form agrees with
the structure of~\cite[Lemma~22]{GaoZhang}.

\begin{lemma}[Exact fourth-order identity]\label{appx-an:lem:identity}
For \(0<\lambda<1\),
\begin{equation}
 P_A'(\lambda)
 =-\frac\lambda3\sum_i\E D^3F(U_\lambda)
                    [\zeta_i,\zeta_i,\zeta_i]
       +\sum_i(\mathcal A_i+\mathcal B_i+\mathcal C_i),
 \label{appx-an:fourth-identity}
\end{equation}
where
\begin{align}
 \mathcal A_i
 &=\int_0^1 k_\lambda(s)
       \E D^4F(H_i+s\gamma\zeta_i)
                 [\zeta_i,\zeta_i,\zeta_i,\zeta_i]\dd s,
 \label{appx-an:A}\\
 k_\lambda(s)
 &=\frac\gamma6(1-s)^2\{(3-\lambda)s-2\lambda\},
 \label{appx-an:kernel}\\
 \mathcal B_i
 &=\frac{1-\lambda}{2}\int_0^1(1-s)
       \E D^4F(H_i+s\gamma\widetilde\zeta_i)
          [\zeta_i,\zeta_i,\widetilde\zeta_i,\widetilde\zeta_i]
       \dd s,
 \label{appx-an:B}\\
 \mathcal C_i
 &=\frac{\lambda(1-\lambda)}{3\gamma}\int_0^1
       \E D^4F(H_i+s\gamma\widetilde\zeta_i)
          [\zeta_i,\zeta_i,\zeta_i,\widetilde\zeta_i]\dd s.
 \label{appx-an:C}
\end{align}
\end{lemma}

\begin{proof}
Expand the segment Hessian in~\eqref{appx-an:signed-hessian}
about \(H_i\), through its linear term. The two moments give
\begin{align*}
 P_A'(\lambda)
 ={}&\frac12\sum_i\E\{(D^2F(U_\lambda)-D^2F(H_i))[\zeta_i^2]\}\\
 &-\frac\lambda3\sum_i\E D^3F(H_i)[\zeta_i^3]
       +\sum_i\mathcal A_i,
\end{align*}
because the second-order integral remainder has kernel
\[
 -\frac{\gamma^2}{2}\int_s^1(t-s)\nu_\lambda(t)\dd t
 =\frac\gamma6(1-s)^2\{(3-\lambda)s-2\lambda\}=k_\lambda(s).
\]
In the Hessian difference, the deleted summand \(Y_i\) is centered
and independent of \(H_i,\zeta_i\). A second-order expansion has
zero expected linear term and gives \(\mathcal B_i\), since
\(\varpi\gamma^2=1-\lambda\). Finally, add and subtract
\(D^3F(U_\lambda)\) in the cubic term. A first-order integral
expansion of that difference gives \(\mathcal C_i\), since
\(\varpi\gamma=(1-\lambda)/\gamma\). These are exactly the three
corrections stated in the lemma.
\end{proof}

All three corrections are evaluated near the same full path:
\begin{equation}
 H_i+s\gamma\zeta_i=U_\lambda+v_i,\qquad
 v_i=s\gamma\zeta_i-Y_i,\qquad \norm{v_i}_U\le4\delta,
 \label{appx-an:translates}
\end{equation}
and the same holds with \(\widetilde\zeta_i\). The bound follows
from \(\gamma\le2\) and the summand face bound. The scalar kernels
in the three corrections have uniformly bounded absolute integrals.
We must preserve the cubic term at its common base until any
aggregate third-moment cancellation has been used.

\subsection{The early bounds}
\label{appx-an:sec:early}

We use one realization of the path for both estimates. With independent
standard Gaussian vectors \(G_0,G_1\), independent of the path array,
write
\begin{equation}
 U_\lambda=V_{0,\lambda}+\sigma_\lambda G_1,\qquad
 V_{0,\lambda}=V_\lambda+\sqrt{\lambda/2}\,G_0,\qquad
 \sigma_\lambda^2=\lambda/2.
 \label{appx-an:gaussian-split}
\end{equation}
This leaves its law unchanged. The bounded and Gaussian components
of \(V_{0,\lambda}\) have equal face variances \(1-\lambda\) and
\(\lambda/2\), respectively. Since their maximum is at least
\(1/3\), Lemma~\ref{lem:concentration} makes \(V_{0,\lambda}\)
a concentrated base for every \(0\le\lambda\le1\).
The bounded component has summand face bound \(2\delta\).
By~\eqref{appx-an:weight-concentration},
\begin{equation}
 \E F(U_\lambda)\theta(U_\lambda)\le Ch\sqrt L.
 \label{appx-an:path-weight}
\end{equation}

The early estimate bounds all derivatives directly. Each derivative
in Lemma~\ref{appx-an:lem:identity} has a repeated pair
\((\zeta_i,\zeta_i)\); its remaining one or two directions have
face norm at most \(\delta\). Formula~\eqref{appx-an:pair-contraction}
controls the pair, and hence, for any face indices \(a,b\),
\[
 \sum_i\E|(u_a\cdot\zeta_i)(u_b\cdot\zeta_i)|\le1.
\]
This gives the moment factor \(\delta^{q-2}\) for derivative order
\(q=3,4\), including the mixed-copy terms.

To separate these moments from the derivatives, first remove the
shifts using positivity of the logistic weights. If
\(\norm v_U\le4\delta\) and \(h\ge4\delta\), every weight
\(F\prod_{a=1}^k\pi_{\alpha_a}\), with \(k\le4\), changes by
at most an absolute factor. Apply the derivative formula
\eqref{appx-eq:partition-derivative} at \(U_\lambda+v_i\) and use
this comparison term by term. The fresh directions are independent
of \(U_\lambda\), so the preceding moment bound applies. Finally,
the sum of weights with \(k\) face factors is
\(F\theta^k\le F\theta\). Consequently, for each of the direction
patterns in the identity,
\[
 \sum_i\E\left|D^qF(U_\lambda+v_i)
       [\zeta_i,\zeta_i,z_{i,3},\ldots,z_{i,q}]\right|
 \le C\delta^{q-2}h^{-q}\E F(U_\lambda)\theta(U_\lambda)
 \le C\delta^{q-2}h^{1-q}\sqrt L.
\]
The cubic term has the additional coefficient \(\lambda/3\).
The bounded integrals of the fourth-order kernels therefore give
\[
 |R_{4,A,h}(\lambda)|\le C\delta^2\sqrt L/h^3,\qquad
 |P_A'(\lambda)|\le C\sqrt L
       \left(\lambda\delta/h^2+\delta^2/h^3\right).
\]
These are~\eqref{eq:remainder-early} and~\eqref{an:early}.
The moment argument is the one in~\cite[Lemma~4.1]{AL}; the
logistic envelope supplies its boundary factor without an eigenvalue
condition.

\subsection{The late bounds}
\label{appx-an:sec:restoration}

For the late estimate we move the repeated pair of derivatives onto
the Gaussian density. This replaces two factors \(h^{-1}\) by
\(\sigma^{-2}\). To keep the bound independent of rank, we then
sum the coefficient matrices using the covariance bound before
taking the nuclear norm of the Gaussian average.

The matrix estimate needed for this step is proved in
Subsection~\ref{appx-an:sec:matrix}. For a concentrated base \(V_0\),
an independent \(G\sim N(0,I_r)\), \(U=V_0+\sigma G\),
\(\delta\sqrt L\le\sigma\le1\), and \(h\ge\delta\), it states
that, for every integer \(k\ge1\),
\begin{equation}
 \E_{V_0}\sum_{\alpha_1,\ldots,\alpha_k}
 \left\|\E_G\left[(GG^\top-I_r)F(U)
              \prod_{a=1}^k\pi_{\alpha_a}(U)\right]\right\|_*
 \le C_1^{k+1}hL^{3/2}.
 \label{appx-an:matrix-bound}
\end{equation}
The sum is over ordered face tuples, \(C_1\ge1\) is absolute,
and \(\norm{\cdot}_*\) denotes the nuclear norm. The next lemma
handles the fact that the derivatives occur at different nearby
bases. Its auxiliary vectors \(Z_i,Z_{i,a}\) are unrelated to the
Gaussian target \(Z\).

\begin{lemma}[Paired derivatives at translated bases]
\label{appx-an:lem:restoration}
Let \(\mathcal H\) be a \(\sigma\)-field, let \(V_0\) be an
\(\mathcal H\)-measurable concentrated base, let \(G\sim N(0,I_r)\) be
independent of \(\mathcal H\), let \(\delta\sqrt L\le\sigma\le1\), and put
\(U=V_0+\sigma G\). For \(s\in\{1,2\}\), let
\((Z_i,Z_{i,1},\ldots,Z_{i,s},v_i)_i\) be random arrays which, conditional
on \(\mathcal H\), are independent of~\(G\). They may depend on
\(\mathcal H\), and on one another. Suppose that, almost surely,
\begin{equation}
 \sum_i\E[Z_iZ_i^\top\mid\mathcal H]\preceq I_r,
 \qquad \norm{Z_{i,a}}_U\le\delta,
 \qquad \norm{v_i}_U\le b\delta.
 \label{appx-an:conditional-budget}
\end{equation}
There is an absolute constant \(\mathsf A_b\), depending only on \(b\),
such that \(h\ge\mathsf A_b\delta\) implies
\begin{equation}
 \left|\sum_i\E D^{s+2}F(U+v_i)
         [Z_i,Z_i,Z_{i,1},\ldots,Z_{i,s}]\right|
 \le\frac{C_b}{\sigma^2}\,
                 \delta^s h^{1-s}L^{3/2}.
 \label{appx-an:restoration-bound}
\end{equation}
\end{lemma}

\begin{proof}
Condition on \(\mathcal H\) and on the arrays. Two Gaussian
integrations by parts give
\begin{align*}
 &\E_G D^{s+2}F(U+v_i)[Z_i,Z_i,Z_{i,1},\ldots,Z_{i,s}]\\
 &\qquad=\sigma^{-2}\E_G\left[
       \tr\{(GG^\top-I_r)Z_iZ_i^\top\}
       D^sF(U+v_i)[Z_{i,1},\ldots,Z_{i,s}]\right].
\end{align*}
Only first and second derivatives remain. With
\(W_j=F\pi_j\) and \(W_{jk}=F\pi_j\pi_k\), they are
\begin{align*}
 DF(x)[z]&=-h^{-1}\sum_jW_j(x)(u_j\cdot z),\\
 D^2F(x)[z,w]&=h^{-2}\left\{
       2\sum_{j,k}W_{jk}(x)(u_j\cdot z)(u_k\cdot w)
       -\sum_jW_j(x)(u_j\cdot z)(u_j\cdot w)\right\}.
\end{align*}
Thus there are only three forms to estimate. For each, write
\(W_\alpha=F\prod_{a=1}^k\pi_{\alpha_a}\), with \(k=1\) or \(2\),
and let \(c_{i,\alpha}\) denote the corresponding product of the
\(s\) remaining face projections. Then \(|c_{i,\alpha}|\le\delta^s\);
the numerical coefficients have absolute value at most two.

We next express every translated weight at the common base. Put
\[
 d_j(v)=e^{u_j\cdot v/h}-1,\qquad
 z_v(x)=\sum_j\pi_j(x)d_j(v),\qquad
 \epsilon=e^{b\delta/h}-1.
\]
Since \(N(x+v)=N(x)(1+z_v(x))\),
\begin{equation}
 W_\alpha(x+v)=W_\alpha(x)
 e^{\sum_a u_{\alpha_a}\cdot v/h}(1+z_v(x))^{-k-1}.
 \label{eq:weight-translation}
\end{equation}
If \(\norm v_U\le b\delta\), then \(|d_j(v)|,|z_v(x)|\le\epsilon\).
For \(\epsilon<1\), expand
\[
 (1+z_v)^{-k-1}
 =\sum_{t\ge0}(-1)^t\binom{k+t}{t}
       \sum_{\beta_1,\ldots,\beta_t}
       \prod_{q=1}^t\pi_{\beta_q}(x)d_{\beta_q}(v).
\]
This exact series leaves every Gaussian-dependent weight at \(U\).
For fixed face tuples \(\alpha,\beta\), all dependence on the
summands is now in the symmetric coefficient matrix
\[
 \Xi_{\alpha,\beta}=\sum_i\E\left[
 Z_iZ_i^\top c_{i,\alpha}
 e^{\sum_a u_{\alpha_a}\cdot v_i/h}
 \prod_{q=1}^t d_{\beta_q}(v_i)\,\middle|\,\mathcal H\right].
\]
For any vector \(z\), the conditional covariance bound implies
\[
 |z^\top\Xi_{\alpha,\beta}z|
 \le\delta^s e^{kb\delta/h}\epsilon^t
       \sum_i\E[(z\cdot Z_i)^2\mid\mathcal H]
 \le\delta^s e^{kb\delta/h}\epsilon^t\norm z^2.
\]
Hence \(\norm{\Xi_{\alpha,\beta}}_\op\) has the same scalar bound.
This is the matrix use of the covariance principle in
\eqref{appx-an:pair-contraction}; replacing it by a trace bound
would introduce the rank.

Apply \(|\tr(\Xi M)|\le\norm\Xi_\op\norm M_*\), sum over face
tuples, and use~\eqref{appx-an:matrix-bound} with \(k+t\) factors.
After restoring \(\sigma^{-2}h^{-s}\), each of the three forms
is bounded by \(\sigma^{-2}\delta^sh^{1-s}L^{3/2}\) times
\[
 e^{kb\delta/h}C_1^{k+1}
 \sum_{t\ge0}\binom{k+t}{t}(C_1\epsilon)^t
 =e^{kb\delta/h}C_1^{k+1}(1-C_1\epsilon)^{-k-1}.
\]
Choose \(\mathsf A_b\) so that \(h\ge\mathsf A_b\delta\) implies
\(C_1\epsilon\le1/2\). Since \(k\le2\), the last expression is
bounded by \(C_b\), proving~\eqref{appx-an:restoration-bound}.
Only the averaged matrix estimate uses concentration of \(V_0\).

The series and expectation interchanges can first be justified by
the convergent scalar majorant \((1-\epsilon)^{-k-1}\), the bound
\(\sum_i\E\norm{Z_i}^2\le r\), and Gaussian integrability.
This finite preliminary bound is used only for the interchanges;
the operator-norm argument supplies the stated bound.
\end{proof}

\begin{proof}[Proof of the late bounds \eqref{an:late}
and~\eqref{an:remainder}]
Use the decomposition~\eqref{appx-an:gaussian-split} and let
\(\mathcal H\) contain the path array and \(G_0\).
The base \(V_{0,\lambda}\) is concentrated, \(G_1\) is independent
of \(\mathcal H\), and \(\sigma_\lambda^2=\lambda/2\).
Thus \(\lambda\ge d_0\delta^2L\), with \(d_0\ge2\), ensures
\(\sigma_\lambda\ge\delta\sqrt L\). The fresh copies remain
independent of \(\mathcal H\), so
\(\sum_i\E[\zeta_i\zeta_i^\top\mid\mathcal H]=I_r\).

For the cubic term, take \(s=1\), \(Z_i=Z_{i,1}=\zeta_i\), and
\(v_i=0\) in Lemma~\ref{appx-an:lem:restoration}. It gives
\begin{equation}
 \left|\frac\lambda3\sum_i\E D^3F(U_\lambda)[\zeta_i^3]\right|
 \le C\lambda\sigma_\lambda^{-2}\delta L^{3/2}
 \le C\delta L^{3/2}.
 \label{appx-an:cubic-bound}
\end{equation}
For each fourth-order correction, use \(s=2\), \(Z_i=\zeta_i\),
and the shifts in~\eqref{appx-an:translates}, with \(b=4\).
The remaining directions are \((\zeta_i,\zeta_i)\) in
\(\mathcal A_i\), \((\widetilde\zeta_i,\widetilde\zeta_i)\) in
\(\mathcal B_i\), and \((\zeta_i,\widetilde\zeta_i)\) in
\(\mathcal C_i\). Their face norms are at most \(\delta\), so each
summed integrand is bounded by
\(C\delta^2L^{3/2}/(\lambda h)\). Integrating the bounded scalar
kernels proves~\eqref{an:remainder}; combining it with the cubic
bound proves~\eqref{an:late}.
\end{proof}

At the endpoints, \(\Gamma U_0\) has the law of \(W\), and
\(\Gamma U_1\) has the law of \(Z\). Boundedness and continuity
of \(F\), together with the integrable early derivative bound,
justify integration up to both endpoints. This completes the proof
of Proposition~\ref{an:prop:derivative}, subject to the matrix
estimate proved next.

\subsection{A matrix estimate for the logistic weights}
\label{appx-an:sec:matrix}

We now prove~\eqref{appx-an:matrix-bound}. The key observation is that
tilting a standard Gaussian law by a log-concave weight can only
decrease its covariance. This reduces the nuclear norm to two scalar
quantities: the squared length of the tilted mean and a signed trace
term. We control the means relative to a single face and sum the
weights before estimating the trace. The latter then becomes a
surface integral, bounded in Subsection~\ref{appx-an:sec:surface}.
This is the soft-weight counterpart of the convex-cell argument
in~\cite[Lemma~12]{GaoZhang}.

Throughout this subsection, \(V_0\) is a concentrated base,
\(G\sim N(0,I_r)\) is independent of \(V_0\),
\(\delta\sqrt L\le\sigma\le1\), and \(U=V_0+\sigma G\).
We use the face scores \(t_j\) from
Subsection~\ref{appx-an:sec:coords}, so
\(N=1+\sum_j e^{t_j}\), \(F=N^{-1}\),
\(\pi_j=e^{t_j}/N\), and \(\theta=1-F\).

For the mean estimate, comparing a weight with the constant one would
give the entropy bound \(\norm{\mu_w}^2\le2\log(1/p_w)\).
Instead, we compare it with a weight depending on one normal coordinate.
This leaves that coordinate to a second-moment bound and measures the
tangential displacement by the loss of mass relative to the reference
weight. This refinement of the entropy argument
in~\cite[Lemma~12]{GaoZhang} avoids a logarithm of the smoothing scale.

\begin{lemma}[Gaussian mean relative to a single face]
\label{appx-an:lem:mean}
Let \(u\) be a unit vector, and let \(b(g)=\beta(u\cdot g)\in[0,1]\) have
positive Gaussian integral \(p_b\). If \(0\le w\le b\),
\(p_w=\E w(G)>0\), and \(\mu_w=\E[Gw(G)]/p_w\), then
\begin{equation}
 p_w\norm{\mu_w}^2
 \le \E[w(G)(u\cdot G)^2]
        +2p_w\log(p_b/p_w).
 \label{appx-an:mean-entropy}
\end{equation}
\end{lemma}

\begin{proof}
For \(t\perp u\), Jensen's inequality, \(w\le b\), and Gaussian
independence give
\[
 p_we^{t\cdot\mu_w}
 \le\E[w(G)e^{t\cdot G}]
 \le\E[b(G)e^{t\cdot G}]
 =p_be^{\norm t^2/2}.
\]
Take \(t\) to be the projection of \(\mu_w\) onto \(u^\perp\).
Taking logarithms bounds its squared length by
\(2\log(p_b/p_w)\). Cauchy--Schwarz gives
\(p_w|u\cdot\mu_w|^2\le\E[w(G)(u\cdot G)^2]\) for the normal
component. Adding the two bounds proves the claim.
\end{proof}

We can now control the matrix sum without a factor depending on the
rank. The trace term must retain its sign until the weights have been
summed; the next subsection provides the needed bound for that sum.

\begin{lemma}[Matrix estimate for logistic weights]
\label{appx-an:lem:matrix}
If \(h\ge\delta\), then~\eqref{appx-an:matrix-bound} holds for every
integer \(k\ge1\).
\end{lemma}

\begin{proof}
We first record the Gaussian-tilt calculation. Let \(w:\R^r\to[0,1]\)
be log-concave with \(p_w=\E_Gw(G)>0\), and let \(\mu_w,K_w\)
be the mean and covariance under the tilted law. Pr\'ekopa's
theorem~\cite[Theorem~6]{Prekopa1973} makes
\[
 z\longmapsto\int_{\R^r}e^{-\norm{g-z}^2/2}w(g)\dd g
\]
log-concave. The Hessian of its logarithm at zero is \(K_w-I_r\),
so \(K_w\preceq I_r\). The same covariance contraction for convex-cell
indicators appears in~\cite[Lemma~11]{GaoZhang}.
Writing
\[
 \E_G[(GG^\top-I_r)w(G)]
 =p_w\{\mu_w\mu_w^\top-(I_r-K_w)\}
\]
as the difference of two positive semidefinite matrices, and using
\(\E_G[H_2(G)w(G)]=p_w(\tr K_w-r+\norm{\mu_w}^2)\), gives
\begin{equation}
 \left\|\E_G[(GG^\top-I_r)w(G)]\right\|_*
 \le2p_w\norm{\mu_w}^2-\E_G[H_2(G)w(G)],
 \qquad H_2(g)=\norm g^2-r.
 \label{appx-an:matrix-trace-reduction}
\end{equation}
Thus the matrix calculation applies to any log-concave weight; only
the two scalar estimates below use the logistic form.

Condition on \(V_0\). For each ordered tuple
\(\alpha=(\alpha_1,\ldots,\alpha_k)\), put
\[
 w_\alpha(g)=F(V_0+\sigma g)
                 \prod_{a=1}^k\pi_{\alpha_a}(V_0+\sigma g),
 \quad p_\alpha=\E_Gw_\alpha(G),
 \quad \mu_\alpha=\frac{\E_G[Gw_\alpha(G)]}{p_\alpha}.
\]
Its logarithm is affine minus \((k+1)\log N(V_0+\sigma g)\), so
\(w_\alpha\) is log-concave. Moreover,
\(\bar w(g):=\sum_\alpha w_\alpha(g)
=F(V_0+\sigma g)\theta(V_0+\sigma g)^k\).
Summing~\eqref{appx-an:matrix-trace-reduction} and then averaging over
\(V_0\) reduces the claim to bounds for
\[
 \E_{V_0}\sum_\alpha p_\alpha\norm{\mu_\alpha}^2
 \quad\text{and}\quad
 \E_{V_0,G}[H_2(G)F(U)\theta(U)^k].
\]

To bound the means, use the first face of each tuple as a reference.
For a tuple with \(\alpha_1=j\), set
\[
 b_j(g)=\frac{e^{t_j(V_0+\sigma g)}}
                  {\{1+e^{t_j(V_0+\sigma g)}\}^2},
 \qquad p_{b_j}=\E_Gb_j(G).
\]
This weight depends only on \(u_j\cdot g\), and
\(w_\alpha\le F\pi_j\le b_j\). Lemma~\ref{appx-an:lem:mean}
and \(|u_j\cdot G|\le\norm G_U\) therefore imply
\begin{equation}
 \sum_\alpha p_\alpha\norm{\mu_\alpha}^2
 \le\E_G[\bar w(G)\norm G_U^2]
       +2\sum_\alpha p_\alpha\log(p_{b_{\alpha_1}}/p_\alpha).
 \label{appx-an:entropy-sum}
\end{equation}
The first term has expectation at most \(ChL^{3/2}\): condition on
\(G\), apply~\eqref{appx-an:weight-concentration} using
\(\bar w\le F\theta\), and then use \(\E\norm G_U^2\le CL\).

For the logarithmic term, define the total averaged masses
\[
 \bar p=\E_{V_0,G}\bar w(G),\qquad
 \bar b=m^{k-1}\sum_j\E_{V_0,G}b_j(G).
\]
There are \(m^{k-1}\) tuples with each first face. The log-sum
inequality and the joint concavity of \((x,y)\mapsto x\log(y/x)\)
therefore give
\[
 \E_{V_0}\sum_\alpha p_\alpha
          \log(p_{b_{\alpha_1}}/p_\alpha)
 \le\bar p\log(\bar b/\bar p).
\]
By~\eqref{appx-an:weight-concentration} and its single-face version,
\(\bar p\le Ch\sqrt L\) and \(\bar b\le Cm^kh\sqrt L\).
Writing \(\bar p=h\sqrt L\,x\), with \(0<x\le C\), yields
\[
 \bar p\log(\bar b/\bar p)
 \le h\sqrt L\,x\log(Cm^k/x)
 \le C(k+1)hL^{3/2}.
\]
The zero-mass case follows by continuity. Together with
\eqref{appx-an:entropy-sum}, this proves
\begin{equation}
 \E_{V_0}\sum_\alpha p_\alpha\norm{\mu_\alpha}^2
 \le C(k+1)hL^{3/2}.
 \label{appx-an:tilted-mean-bound}
\end{equation}
The reference masses have the same factor \(h\sqrt L\) as the
weights being bounded, so this factor cancels inside the logarithm.

For the trace, shifting every threshold by \(s\) gives functions
\(F_s,\theta_s\) satisfying
\[
 F\theta^k=-\frac hk\left.\partial_s\theta_s^k\right|_{s=0}.
\]
We turn this derivative into a surface integral by expressing powers
of \(F_s\) as mixtures of polyhedral indicators. For an integer
\(t\ge1\), let \(T_t\sim\operatorname{Gamma}(t,1)\) and let
\(\omega_j\) be independent unit exponentials, all independent of
\(V_0,G\). The Gamma Laplace transform gives
\begin{equation}
 F_s(x)^t=\Pp\{u_j\cdot x\le
        a_j+s+h\log(\omega_j/T_t),\ j\le m\}.
 \label{appx-an:gamma-mixture}
\end{equation}
Expand \(\theta_s^k=\sum_{t=0}^k(-1)^t\binom kt F_s^t\).
The absolute coefficients sum to \(2^k\), and the constant term has
zero derivative. For fixed random thresholds, the first-variation formula says that
increasing every unit-normal threshold at unit speed differentiates
the Gaussian integral into the boundary integral, with the factor
\(\sigma^{-1}\) after rescaling. This holds for almost every common
shift. Uniform bounds on Gaussian integrals over the finitely many
supporting hyperplanes justify differentiation under the mixture
average. The independent, absolutely continuous shift \(-h\log T_t\)
avoids the exceptional levels, including those at which a target
polyhedron is lower dimensional. Lemma~\ref{appx-an:lem:surface}
therefore gives
\begin{align*}
 \left|\E_{V_0,G}[H_2(G)F(U)\theta(U)^k]\right|
 &\le\frac{2^kh}{k\sigma}
       \sup_Q\E_{V_0}|\surf((Q-V_0)/\sigma)|\\
 &\le\frac{C2^kh}{k\sigma}
       (\sigma L^{3/2}+\delta L^2)
 \le C2^khL^{3/2},
\end{align*}
where the last inequality uses \(\sigma\ge\delta\sqrt L\).
Combining this trace bound with
\eqref{appx-an:matrix-trace-reduction} and
\eqref{appx-an:tilted-mean-bound} proves
\eqref{appx-an:matrix-bound}, after absorbing \(k+1\) and \(2^k\)
into \(C_1^{k+1}\).
\end{proof}

\subsection{An averaged Gaussian surface estimate}
\label{appx-an:sec:surface}

The trace term in Lemma~\ref{appx-an:lem:matrix} becomes a Gaussian
surface integral over a randomly translated polyhedron. We now show
that averaging over the concentrated translation gains the factor
\(\sigma\) needed to cancel the \(\sigma^{-1}\) in that argument.
We retain the notation
\[
 H_2(g)=\norm g^2-r,\qquad
 \surf(P)=\int_{\partial P}H_2(g)\varphi_r(g)\dd\mathcal H^{r-1}(g),
\]
where \(\varphi_r\) is the standard Gaussian density and each facet is
counted once. The integral is defined for proper, full-dimensional
polyhedra: those with nonempty interior that are not the whole space.
We set \(\surf(P)=0\) for the empty set and the whole space. As in
Subsection~\ref{appx-an:sec:coords}, \(P^t\) means that all thresholds
in the defining system of inequalities with unit face normals are increased by
\(t\); it does not denote a Euclidean neighborhood.

\begin{lemma}[Averaged scalar surface]\label{appx-an:lem:surface}
Let \(V_0\) be a concentrated base, and let \(0<\sigma\le1\). Uniformly
over target polyhedra \(Q\) with face normals among \(u_1,\ldots,u_m\), empty
or full dimensional,
\begin{equation}
 \E\left|\surf\bigl((Q-V_0)/\sigma\bigr)\right|
 \le C\bigl(\sigma L^{3/2}+\delta L^2\bigr).
 \label{appx-an:surface-bound}
\end{equation}
In particular, if \(\sigma\ge\delta\sqrt L\), the right side is at most
\(C\sigma L^{3/2}\).
\end{lemma}

\begin{proof}
We first bound the surface integral of a fixed polyhedron by an integral
over short outward tubes, and then average over \(V_0\). The geometric
argument uses the facet--ridge divergence identity and normal tubes
from the proof of~\cite[Proposition~4]{GaoZhang}. We choose the tube
width at each boundary point, so that the subsequent averaging reduces
to two Gaussian moments.

Only projections onto face normals and planes spanned by pairs of face normals
will be needed. For linearly independent \(u_j,u_k\), let \(\Pi_{jk}\)
be orthogonal projection onto \(\operatorname{span}(u_j,u_k)\), and put
\begin{equation}
 \mathcal M(g)=\max\left\{\max_j|u_j\cdot g|,
       \max_{j,k:\,u_j,u_k\text{ independent}}\norm{\Pi_{jk}g}\right\},
 \qquad R(g)=1+\mathcal M(g).
 \label{appx-an:projection-radius}
\end{equation}
An empty second maximum is omitted. The function \(R\) is
\(1\)-Lipschitz, since each projection is a contraction. For
\(G\sim N(0,I_r)\), the maximum involves at most \(m+m^2\)
Gaussian projections of dimension at most two. Their tail bounds and
a union bound give
\[
 \Pp\{\mathcal M(G)>t\}\le\min\{1,Cm^2e^{-ct^2}\}.
\]
Integrating this bound, with the integral split at a constant times
\(\sqrt{\log(em)}\), and using \(m\le2p\), yields
\begin{equation}
 \E R(G)^2\le CL,\qquad \E R(G)^3\le CL^{3/2}.
 \label{appx-an:projection-moments}
\end{equation}

For every proper, full-dimensional polyhedron \(P\) with face normals among
the \(u_j\), we will prove
\begin{equation}
 |\surf(P)|\le C\E_G\left[R(G)^3
       \ind\{G\in P^{1/R(G)}\setminus P\}\right].
 \label{appx-an:adaptive-tubes}
\end{equation}
The right side integrates over an outer layer whose width is smaller
where the normal projections are larger. This keeps the Gaussian
density comparable along each tube without using the full norm
\(\norm G\), whose moments would introduce the rank \(r\).

To prove~\eqref{appx-an:adaptive-tubes}, describe the actual facets of
\(P\) by deleting repeated and redundant inequalities. Write \(F_j\)
for the facet in \(u_j\cdot x=a'_j\), and let
\(s_j=\int_{F_j}\varphi_r\dd\mathcal H^{r-1}\).
A ridge \(E_{jk}\) is a face of dimension \(r-2\) shared by exactly
two facets, \(F_j\) and \(F_k\). Let \(\alpha_{jk}\in(0,\pi)\) be the
angle between their face normals, and set
\[
 \rho_{jk}=\int_{E_{jk}}\varphi_r\dd\mathcal H^{r-2},\qquad
 z_{jk}=\Pi_{jk}x\quad(x\in E_{jk}).
\]
The vector \(z_{jk}\) is constant on the ridge: it is determined by
\(u_j\cdot z_{jk}=a'_j\) and \(u_k\cdot z_{jk}=a'_k\).
The divergence identity is
\begin{equation}
 \surf(P)=\sum_j\bigl((a'_j)^2-1\bigr)s_j
       -\sum_{\{j,k\}}(d_{k\mid j}+d_{j\mid k})\rho_{jk},
 \label{appx-an:surface-identity}
\end{equation}
where the second sum runs over the ridges and
\[
 d_{k\mid j}=\frac{a'_k-(u_j\cdot u_k)a'_j}{\sin\alpha_{jk}}.
\]
Indeed, writing \(x=a'_ju_j+y\) on facet \(j\), with \(y\perp u_j\),
splits \(H_2(x)\) into \((a'_j)^2-1\) and \(\norm y^2-(r-1)\).
The first term gives the facet contribution above. For the second,
\[
 (\norm y^2-(r-1))\varphi_{r-1}(y)
 =-\operatorname{div}_{u_j^\perp}\{y\varphi_{r-1}(y)\}.
\]
The divergence theorem within the facet gives
\(-d_{k\mid j}\rho_{jk}\) on its ridge with facet \(k\), because
the outward unit normal to that ridge within the facet is
\((u_k-(u_j\cdot u_k)u_j)/\sin\alpha_{jk}\).
Summing over facets proves~\eqref{appx-an:surface-identity}.
Truncation handles unbounded facets, since polynomial Gaussian terms
vanish at infinity. For \(r=1\), there are no ridges and the identity
is a sum over endpoints.

We pair the two contributions of each ridge before bounding them.
The elementary identity
\[
 d_{k\mid j}+d_{j\mid k}
 =2\sin(\alpha_{jk}/2)
       \frac{u_j+u_k}{\norm{u_j+u_k}}\cdot z_{jk}
\]
gives
\begin{equation}
 |d_{k\mid j}+d_{j\mid k}|\le\alpha_{jk}\norm{z_{jk}}.
 \label{appx-an:paired-ridge-bound}
\end{equation}
Thus the inverse sine in the separate terms disappears. This pairing
allows the estimate to remain uniform when facets are nearly parallel.

At a point \(x\) in the relative interior of a facet or ridge, choose
the tube width
\[
 \eta(x)=\frac1{4R(x)}.
\]
Over a facet, use displacements \(y=tu_j\), with \(0<t<\eta(x)\).
Over a ridge, use \(0<\norm y<\eta(x)\) in the outward normal cone,
the sector of angle \(\alpha_{jk}\) spanned by \(u_j,u_k\).
These are precisely normal directions for which \(x\) is the nearest
point in \(P\) to \(x+y\). In both cases,
\[
 |x\cdot y|\le\mathcal M(x)\norm y\le\tfrac14,\qquad
 \frac{\varphi_r(x+y)}{\varphi_r(x)}
       =e^{-x\cdot y-\norm y^2/2}\ge e^{-1/4-1/32}.
\]
For a ridge, the first bound uses \(x\cdot y=z_{jk}\cdot y\).
Since \(R\) is \(1\)-Lipschitz and \(R\ge1\), we also have
\(3R(x)/4\le R(x+y)\le5R(x)/4\).
Integrating the weight \(R^3\varphi_r\) along a facet's normal
interval therefore gives
\[
 \int_0^{\eta(x)} R(x+tu_j)^3\varphi_r(x+tu_j)\dd t
 \ge cR(x)^2\varphi_r(x)
 \ge c\{1+(a'_j)^2\}\varphi_r(x).
\]
The corresponding integral over a ridge's normal sector gives
\[
 \int_{\substack{y\in N_P(x)\\0<\norm y<\eta(x)}}
        R(x+y)^3\varphi_r(x+y)\dd y
 \ge c\alpha_{jk}R(x)\varphi_r(x)
 \ge c\alpha_{jk}\norm{z_{jk}}\varphi_r(x),
\]
where \(N_P(x)\) is the normal cone and \(\dd y\) is area in its
two-dimensional plane. The first integral contains one power of
\(\eta(x)\); the second contains two. This explains the common weight
\(R^3\) in~\eqref{appx-an:adaptive-tubes}.

The tubes over distinct faces are disjoint up to null sets, by
uniqueness of the nearest point in a closed convex set. Within each
face, the decomposition into tangential and normal coordinates has
Jacobian one. Fubini's theorem therefore applies even though
\(\eta(x)\) varies along the face. Every tube point \(g=x+y\) lies
outside \(P\), and, for every inequality in its original description,
\[
 u_j\cdot g-a'_j\le\norm y<\frac1{4R(x)}\le\frac1{R(g)}.
\]
Thus all tubes lie in \(\{g:g\in P^{1/R(g)}\setminus P\}\), including
when the defining system has redundant inequalities. Integrating the
two tube estimates over the facets and ridges, and applying
\eqref{appx-an:surface-identity} and
\eqref{appx-an:paired-ridge-bound}, proves
\eqref{appx-an:adaptive-tubes}.

We now apply this deterministic bound to \(P=(Q-V_0)/\sigma\),
with \(G\) independent of \(V_0\). Empty targets and the whole space
contribute zero, so assume \(Q\) is proper and full dimensional.
Rescaling the boundary event gives
\[
 G\in P^{1/R(G)}\setminus P
 \quad\Longleftrightarrow\quad
 V_0+\sigma G\in Q^{\sigma/R(G)}\setminus Q.
\]
Conditional on \(G\), the width \(\sigma/R(G)\) is fixed. Translate
the thresholds by \(-\sigma G\) and enlarge this outer layer to the
symmetric layer in~\eqref{appx-an:concentration}. Concentration then
bounds its conditional probability by
\(C\{\sigma/R(G)+\delta\}\sqrt L\). Tonelli's theorem and
\eqref{appx-an:projection-moments} yield
\begin{align*}
 \E_{V_0}|\surf((Q-V_0)/\sigma)|
 &\le C\E_G\left[R(G)^3
       \Pp_{V_0}\{V_0+\sigma G\in Q^{\sigma/R(G)}\setminus Q\}\right]\\
 &\le C\sqrt L\{\sigma\E R(G)^2+\delta\E R(G)^3\}\\
 &\le C\{\sigma L^{3/2}+\delta L^2\}.
\end{align*}
This is~\eqref{appx-an:surface-bound}.
\end{proof}

% Source: concentration_proof
\section{\texorpdfstring{Proof of the anti-concentration inequality of Lemma~\ref{lem:concentration}}{Proof of the anti-concentration inequality of Lemma 1}}\label{appx-app:concentration}

\label{appx-conc:conditioning}

We prove Lemma~\ref{lem:concentration} by representing each summand
as a random choice between two independent copies. Conditional on these
copies, the sum becomes a translated Rademacher sum, to which we can
apply the anti-concentration inequality of O'Donnell, Servedio, and
Tan~\cite{OST} in the form recorded in~\cite[Lemma~J.4]{CCKK2022}.
This gives a short proof of the bounded-sum estimate stated
in~\cite[Lemma~2.1]{AL}.

First take \(v=1\) and omit the independent shift. Let
\(\zeta_i'\) be independent copies of the summands, and let
\(\varepsilon_i\) be independent uniform signs, independent of all
these vectors. Put
\[
 \bar\zeta_i=\frac{\zeta_i+\zeta_i'}2,\qquad
 \Delta_i=\frac{\zeta_i-\zeta_i'}2,\qquad
 \mathcal G=\sigma(\zeta_i,\zeta_i':1\le i\le n).
\]
Each \(\bar\zeta_i+\varepsilon_i\Delta_i\) selects one of the two copies, so
\begin{equation}
 S\stackrel d=\widetilde S:=\sum_i\bar\zeta_i+\sum_i\varepsilon_i\Delta_i.
 \label{appx-eq:pair-representation}
\end{equation}
Conditional on \(\mathcal G\), the first sum is a fixed translation,
\(\norm{\Delta_i}_\infty\le\eta\), and the coordinate variances of the
second sum are \(V_j=\sum_i\Delta_{ij}^2\). We need these variances to
stay away from zero. Independence and centering give
\[
 \E V_j=\frac12,\qquad
 0\le\Delta_{ij}^2\le\eta^2,\qquad
 \sum_i\Var(\Delta_{ij}^2)\le\frac{\eta^2}{2}.
\]
Bernstein's inequality and a union bound therefore imply, for an
absolute constant \(c_1>0\),
\begin{equation}
 \Pp(\mathcal E^c)\le p\exp(-c_1/\eta^2),
 \qquad \mathcal E=\{\min_{j\le p}V_j\ge1/4\}.
 \label{appx-eq:conditional-variance}
\end{equation}

On \(\mathcal E\), apply~\cite[Lemma~J.4]{CCKK2022} conditionally to
\(\sum_i\varepsilon_i\Delta_i\). Encode the lower and upper coordinate
bounds by the \(2p\) coordinates of \((x,-x)\), and absorb the fixed
translation \(\sum_i\bar\zeta_i\) into the thresholds. The outer and inner
rectangles correspond to thresholds separated by \(2t\), so
\[
 \Pp(\widetilde S\in A^{t}\setminus A^{-t}\mid\mathcal G)
 \le C(t+\eta)\sqrt L\qquad\text{on }\mathcal E.
\]
Infinite endpoints follow by taking limits. This argument also covers
inconsistent coordinate bounds, since the orthant inequality is uniform
in all thresholds. Averaging and using
\eqref{appx-eq:conditional-variance} gives
\begin{equation}
 \Pp(S\in A^{t}\setminus A^{-t})
 \le C(t+\eta)\sqrt L+p\exp(-c_1/\eta^2).
 \label{appx-eq:concentration-pairs}
\end{equation}
If \(\eta\sqrt L\le c_2\), with \(c_2>0\) sufficiently small, then
\(L\le c_1/(2\eta^2)\), and the last term is at most
\(\exp(-c_1/(2\eta^2))\le C\eta\le C\eta\sqrt L\).
Otherwise the desired estimate follows from the bound by one.
This proves the bounded-sum assertion for \(v=1\). Dividing by
\(\sqrt v\) gives the general case, and conditioning on \(R\)
permits any independent shift.

For a centered Gaussian vector \(G\), Nazarov's anti-concentration
inequality in
\cite[Lemma~A.1]{CCK2017}, also recorded in
\cite[Lemma~J.3]{CCKK2022}, applied to \((G,-G)/\sqrt v\) gives
\[
 \Pp(G\in A^{t}\setminus A^{-t})\le Ct\sqrt{L/v}.
\]
It requires only positive marginal variances and therefore permits
singular covariance. Conditioning on the independent shift and taking
the minimum with one completes the proof.

% Source: subexponential_extension_proof
\section{\texorpdfstring{Sub-exponential coordinates: proof of Corollary~\ref{cor:subexponential}}{Sub-exponential coordinates: proof of Corollary 2}}\label{app:subexponential}
\begin{proof}[Proof of Corollary~\ref{cor:subexponential}]
We clip each coordinate, apply the bounded theorem after marginal
standardization, and transfer the result back to \(W_n\). This requires
control of both the discarded sum and the Gaussian covariance change.
We drop the subscript \(n\) and write \(\ell\) and \(b\) for the
quantities \(\ell_n\) and \(b_n\) of the corollary.

We first prove the bounds under \(\E e^{|X_{ij}|/K}\le2\) for every
coordinate. The stated norm condition gives this exponential-moment
bound at scale \(2K\): the defining infimum supplies an admissible
scale below \(2K\) for each coordinate, and increasing the scale
decreases its exponential moment. Replacing \(K\) by \(2K\) doubles
\(b\), while the positive logarithms in both asserted rates decrease.
Each rate therefore increases by at most a factor two.

The unrestricted bound follows from the bound by one, after adjusting its absolute
constant, when \(b\ge1/16\), so assume \(b<1/16\).

\emph{Step 1: Clip the summands and bound the residual sum.}
With the clipping level \(4K\ell\), define coordinatewise
\[
 C_T(x)=\operatorname{sgn}(x)\min\{|x|,4K\ell\},\qquad
 Y_i=C_T(X_i)-\E C_T(X_i),\qquad D_i=X_i-C_T(X_i).
\]
Here \(Y_i\) is the centered clipped summand and \(D_i\) is the discarded
part. Their normalized sums, marked by the subscripts \(T\) and \(D\), satisfy
\[
 W_T=n^{-1/2}\sum_iY_i,\qquad
 W_D=n^{-1/2}\sum_i(D_i-\E D_i),\qquad W_n=W_T+W_D.
\]
The exponential-moment assumption gives, for \(t\ge0\),
\[
 \Pp(|D_{ij}|>t)\le2e^{-4\ell-t/K},\qquad
 \E|D_{ij}|^k\le2k!K^ke^{-4\ell}\quad(k\ge1).
\]
In particular, independence of the summands implies
\begin{equation}
 \max_j\Var(W_{D,j})\le4K^2e^{-4\ell}.
 \label{tail:residual-moments}
\end{equation}
The pointwise bound \(e^{|D_{ij}|/K}\le1+e^{-4\ell}e^{|X_{ij}|/K}\)
gives \(\E e^{|D_{ij}|/K}\le1+2e^{-4\ell}\). Together with
\(\E|D_{ij}|\le2Ke^{-4\ell}\), independence and the triangle inequality
therefore yield
\begin{align*}
 \E\exp\!\left(\frac{\sqrt n\,|W_{D,j}|}{K}\right)
 &\le\exp\!\left(\frac1K\sum_i\E|D_{ij}|\right)
          \prod_i\E e^{|D_{ij}|/K}\\
 &\le e^{2ne^{-4\ell}}(1+2e^{-4\ell})^n
 \le e^{4ne^{-4\ell}}.
\end{align*}
For \(\bar L=\log(4np)\), Markov's inequality gives
\[
 \Pp\!\left(\sqrt n\,|W_{D,j}|\ge K[\bar L+4ne^{-4\ell}]\right)\le e^{-\bar L}.
\]
Since \(n,p\ge1\),
\[
 4ne^{-4\ell}=4e^{-4}n^{-3}\le\frac14\le\log4\le \bar L.
\]
A coordinate union bound now yields
\begin{equation}
 \Pp(\|W_D\|_\infty>\varrho_1)\le\frac1{4n}\le\frac1{2n},\qquad
 \varrho_1=\frac{2K\bar L}{\sqrt n}.
 \label{tail:sum-coupling}
\end{equation}
Since the union bound is applied to the residual sums, the clipping
level can depend only on \(K\) and \(\log n\).

\emph{Step 2: Apply the bounded comparison.}
Set \(\tau=2Ke^{-2\ell}\), and write \(\operatorname{sd}(W_{T,j})\) for the
standard deviation of the \(j\)th coordinate of~\(W_T\).
Since each coordinate of \(W_n\) has variance one, the triangle
inequality in \(L^2\) gives
\[
 |\operatorname{sd}(W_{T,j})-1|\le\tau,
 \qquad \frac\tau b=\frac{2\sqrt n\,e^{-2\ell}}\ell
 \le\frac2{e^2}.
\]
Hence \(\operatorname{sd}(W_{T,j})\ge1/2\). Since \(\|Y_i\|_\infty\le8K\ell\), the
independent centered summands of the marginally standardized \(W_T\)
have envelope at most
\[
 \frac{8K\ell}{\sqrt n\min_j\operatorname{sd}(W_{T,j})}\le16b.
\]
Let \(Z_T\) be Gaussian with covariance \(\Cov(W_T)\).
Positive diagonal scaling preserves rectangle distance. Applying
Theorem~\ref{thm:unrestricted} with envelope \(16b\) gives
\begin{equation}
 \rh(W_T,Z_T)\le C\min\left\{1,
 bL^{3/2}\sqrt{1+\log_+\frac1{b\sqrt L}}\right\}.
 \label{tail:bounded-comparison}
\end{equation}
The factor 16 is absorbed into the absolute constant, since increasing
the envelope decreases the logarithmic factor.

\emph{Step 3: Restore the covariance.}
Construct a jointly centered Gaussian vector \((Z_T,Z_D)\) with the
block covariance of \((W_T,W_D)\). This covariance is positive
semidefinite, even when singular. Set \(Z_n=Z_T+Z_D\) on this coupling.
Then \(Z_n\) has exactly the covariance of \(W_n\), while
\(\max_j\Var(Z_{D,j})\le\tau^2\). Gaussian tails therefore give
\begin{equation}
 \Pp(\|Z_D\|_\infty>\varrho_2)\le\frac1{2n},\qquad
 \varrho_2=\sqrt8\,Ke^{-2\ell}\sqrt{\bar L}.
 \label{tail:Gaussian-coupling}
\end{equation}
The Gaussian components need not be independent.

For any rectangle \(A\), the two coupling bounds give
\begin{align*}
 \Pp(W_n\in A)
 &\le\Pp(W_T\in A^{\varrho_1})+\frac1{2n}\\
 &\le\Pp(Z_T\in A^{\varrho_1})+\rh(W_T,Z_T)+\frac1{2n}\\
 &\le\Pp(Z_n\in A^{\varrho_1+\varrho_2})+\rh(W_T,Z_T)+\frac1n.
\end{align*}
Inward expansions give the reverse bound
\(\Pp(W_n\in A)\ge\Pp(Z_n\in A^{-(\varrho_1+\varrho_2)})
-\rh(W_T,Z_T)-1/n\). Applying Gaussian rectangle concentration at
the unit-marginal target \(Z_n\), from Lemma~\ref{lem:concentration}, yields
\begin{equation}
 \rh(W_n,Z_n)\le\rh(W_T,Z_T)
       +C(\varrho_1+\varrho_2)\sqrt L+\frac1n.
 \label{tail:transfer}
\end{equation}
No independence between the two vectors within either coupling is
required for this sandwich argument.
The transfer error has order \(bL^{3/2}\): since
\(\bar L\le2\ell L\) and \(K\ge1\),
\[
 (\varrho_1+\varrho_2)\sqrt L+\frac1n
 \le CK\left[e^{-2\ell}\sqrt{L\bar L}
                     +\frac{\bar L\sqrt L}{\sqrt n}\right]+\frac1n
 \le CbL^{3/2}.
\]
Combine this with \eqref{tail:bounded-comparison} and the bound by one
to prove \eqref{eq:subexponential}.

\emph{Step 4: Preserve a fixed correlation gap.}
For \eqref{eq:subexponential-gap}, it suffices to consider
\(b<\min\{1/16,c_0/4\}\); otherwise its right-hand side is at
least one after increasing \(C_{c_0}\).
Under the Gaussian coupling, both \(Z_n\) and the standardized clipped
vector \((Z_{T,j}/\operatorname{sd}(W_{T,j}))_{j\le p}\) have unit coordinate variances, and
\[
 \|Z_{T,j}/\operatorname{sd}(W_{T,j})-Z_{n,j}\|_{L^2}\le2\tau.
\]
Writing the difference of two covariance entries as a sum of two
products and applying Cauchy--Schwarz shows that the covariance matrices
differ entrywise by at most \(4\tau\). Since
\(4\tau\le(8/e^2)b<c_0/2\), the standardized clipped covariance
has correlation gap at least \(c_0/2\).
Theorem~\ref{thm:optimized} applied with envelope \(16b\) now gives
\[
 \rh(W_T,Z_T)\le C_{c_0}\min\left\{1,
 b\left[L^{3/2}+L\sqrt{1+\log_+\frac1{bL}}\right]\right\}.
\]
The transfer error in \eqref{tail:transfer} is bounded by the
\(CbL^{3/2}\) term. This proves \eqref{eq:subexponential-gap}.
\end{proof}

Coordinatewise odd clipping preserves central symmetry of each summand,
but need not preserve cancellation of third-moment tensors across summands.
Moreover the transfer cost above contains \(bL^{3/2}\), so the
zero-skewness refinement cannot simply be read off by deleting the
cubic term from the displayed sub-exponential bound.

\paragraph{Comparison with the locally nondegenerate result.}
The companion paper~\cite[Corollary~1]{Local} assumes the same
coordinatewise \(\psi_1\) condition, an averaged fourth-moment bound
\(\max_jn^{-1}\sum_i\E|X_{ij}|^4\le B_4^2\), and uniformly positive
eigenvalues for principal covariance submatrices of order at most four.
When these principal-submatrix eigenvalues are bounded below by \(a>0\), its bound is
\[
 \frac{C_a}{\sqrt n}\left[B_4L^{3/2}
                   +K\log(enp)\sqrt L\right],
\]
capped at one. For fixed \(K,B_4,a\) and \(p\) growing polynomially in~\(n\),
this is \(O(n^{-1/2}(\log n)^{3/2})\), whereas
\eqref{eq:subexponential} is \(O(n^{-1/2}(\log n)^3)\).
The latter requires neither a fourth-moment hypothesis in addition to
the tail bound nor a local covariance assumption. If \(K\) varies,
the elementary fourth-moment consequence is \(B_4\le C K^2\);
the tail bound alone does not justify taking \(B_4=K\).

\end{document}

